\documentclass[11pt,openany,twoside]{memoir}

% 6x9 trim size
\usepackage[paperwidth=6in,paperheight=9in,
            inner=0.75in,outer=0.625in,
            top=0.75in,bottom=0.75in]{geometry}

% Fonts (XeLaTeX)
\usepackage{fontspec}
\setmainfont{TeX Gyre Pagella}
\setsansfont{TeX Gyre Heros}
\setmonofont[Scale=0.85]{DejaVu Sans Mono}

% Packages
\usepackage{listings}
\usepackage{booktabs}
\usepackage{longtable}
\usepackage{graphicx}
\usepackage{xcolor}
\usepackage{hyperref}
\usepackage{makeidx}
\usepackage{microtype}
\usepackage{enumitem}
\usepackage{fancyvrb}
\usepackage{titlesec}

% Colors
\definecolor{acmeyellow}{RGB}{234,255,198}
\definecolor{acmeblue}{RGB}{234,234,255}
\definecolor{codeblue}{RGB}{0,0,128}
\definecolor{codegray}{RGB}{96,96,96}

% Listings setup
\lstset{
  basicstyle=\small\ttfamily,
  keywordstyle=\color{codeblue}\bfseries,
  commentstyle=\color{codegray}\itshape,
  breaklines=true,
  frame=single,
  framerule=0.4pt,
  xleftmargin=1em,
  xrightmargin=0.5em,
  aboveskip=0.8em,
  belowskip=0.8em,
  columns=fullflexible,
  keepspaces=true,
  showstringspaces=false,
}

% Hyperref setup
\hypersetup{
  colorlinks=true,
  linkcolor=black,
  urlcolor=codeblue,
  citecolor=black,
  bookmarksdepth=3,
  pdfauthor={},
  pdftitle={The Acme Editor: A Comprehensive Guide},
}

% Index
\makeindex

% Chapter and section styling
\chapterstyle{veelo}
\setsecnumdepth{subsection}

% Custom commands
\newcommand{\button}[1]{\textsf{B#1}}
\newcommand{\key}[1]{\textsf{#1}}
\renewcommand{\cmd}[1]{\texttt{#1}}
\newcommand{\menu}[1]{\textsf{#1}}
\newcommand{\file}[1]{\texttt{#1}}
\newcommand{\env}[1]{\texttt{\$#1}}
\newcommand{\chord}[2]{\textsf{B#1--B#2}}

% Margin notes for tips
\newcommand{\tip}[1]{\marginpar{\footnotesize\raggedright\itshape #1}}

% Running headers
\makepagestyle{acmestyle}
\makeevenhead{acmestyle}{\thepage}{}{\leftmark}
\makeoddhead{acmestyle}{\rightmark}{}{\thepage}
\makeevenfoot{acmestyle}{}{}{}
\makeoddfoot{acmestyle}{}{}{}
\pagestyle{acmestyle}

\begin{document}
\frontmatter

% ---- Title Page ----
\thispagestyle{empty}
\vspace*{2in}
\begin{center}
{\Huge\bfseries The Acme Editor}\\[0.5em]
{\Large A Comprehensive Guide}\\[2em]
{\large Based on the standalone acme derived from Plan~9}\\[4em]
\end{center}
\vfill
\newpage

% ---- Copyright (verso of title page) ----
\thispagestyle{empty}
\vspace*{\fill}
\noindent
The acme editor was created by Rob Pike at Bell Labs as part of
the Plan~9 operating system.  The standalone version documented here
is derived from plan9port (Plan~9 from User Space) by Russ Cox.\\[1em]
\noindent
Acme source code: MIT License.\\
Copyright \textcopyright\ 2021 Plan~9 Foundation.\\
Portions copyright \textcopyright\ 2001--2008 Russ Cox.\\
Portions copyright \textcopyright\ 2008--2009 Google Inc.\\[2em]
\newpage

% ---- Table of Contents ----
\tableofcontents
\newpage

% ---- Chapters ----
\mainmatter
\chapter{Introduction}
\label{ch:intro}

\index{acme!introduction}

Acme is a text editor and programming environment created by Rob Pike
at Bell Labs in 1993.  It was designed as part of the Plan~9 operating
system, where it served as both an editor and a shell---a surface on
which all work takes place.

Acme is unlike any other editor you have used.  It has no menus, no
toolbars, no configuration files, no plugin system, and no key bindings
to memorize.  Instead, it offers a blank surface where all commands are
ordinary text that you execute with a mouse click.  Any text, anywhere
on screen, can be a command.  Any text can be a filename.  Any text can
be a search term.  The entire interface is built from three mouse
buttons and text.

\section{Philosophy}

Rob Pike described acme's design philosophy in his 1994 paper:

\begin{quote}
The right model is to have a shell that does not distinguish between a
program that is an editor and a program that is a compiler.  All programs
should have equal access to the services of the environment.
\end{quote}

\index{philosophy}
This leads to several principles:

\begin{description}[style=nextline]
\item[Text is the universal interface.]
Commands are text.  Output is text.  Filenames are text.  You interact
with everything by reading and clicking on text.

\item[The mouse is primary.]
Acme is designed for a three-button mouse.  The keyboard is for typing
text; the mouse is for everything else---executing commands, opening
files, selecting, searching, scrolling, and window management.

\item[Simplicity over features.]
Acme has no syntax highlighting, no auto-completion, no bracket matching,
no built-in debugger, no project manager.  It provides a powerful surface
and lets external programs do the work.

\item[Compose, don't configure.]
Rather than configuring the editor, you compose it with external tools.
Need to sort lines?  Select them and execute \cmd{|sort}.  Need to
reformat?  Execute \cmd{|fmt}.  The Unix toolkit is your extension
mechanism.
\end{description}

\section{How Acme Differs}

If you come from \textbf{vi} or \textbf{Vim}, you will notice that acme
has no modes.  There is no insert mode, no command mode, no visual mode.
You are always editing.  Commands are not key sequences---they are text
you click on.

If you come from \textbf{Emacs}, you will notice that there are no key
bindings to learn beyond basic cursor movement.  There is no Lisp
configuration, no major modes, no minor modes.  The mouse does what
key chords do in Emacs.

If you come from \textbf{VS Code} or similar IDEs, you will notice the
absence of everything: no file tree, no tabs, no status bar, no
integrated terminal widget.  Instead, every window is a general-purpose
text buffer.  A directory listing is a window.  A shell session is a
window.  Compiler output is a window.  They are all the same thing.

The learning curve is not in memorizing commands---there are very few.
The learning curve is in \emph{changing how you think} about interacting
with a computer.

\section{History}

\index{Plan 9}
\index{Pike, Rob}
\index{Bell Labs}
Acme was written by Rob Pike at Bell Labs as part of the Plan~9
operating system.  Plan~9 was the successor to Unix, designed by many
of the same people.  In Plan~9, all resources---files, network
connections, processes, graphics---appear as files in a hierarchical
namespace.  Acme takes full advantage of this: it presents its windows,
buffers, and events as a file system that other programs can read and
write.

\index{plan9port}
\index{Cox, Russ}
In 2003, Russ Cox ported the Plan~9 user-space tools to Unix systems
as \emph{plan9port} (Plan~9 from User Space).  This made acme available
on Linux, macOS, and other Unix systems, though it still depended on
several Plan~9 support programs (devdraw, fontsrv, plumber, 9pserve).

The standalone version documented in this book goes further: it embeds
all those support programs directly into a single binary.  It uses SDL3
for graphics (supporting both X11 and Wayland on Linux, and native
windows on macOS and Windows), FreeType and fontconfig for fonts, and a
built-in plumber for opening files and URLs.  The result is a
single-binary acme with no external dependencies beyond standard system
libraries.

\section{About This Book}

This book is a comprehensive guide to using acme.  It covers
installation on Linux, macOS, and Windows; the mouse and keyboard
interface; all built-in commands; the sam editing language; plumbing;
and practical workflows.

Chapter~\ref{ch:mouse} (The Mouse) is the most important chapter in
this book.  If you read only one chapter, read that one.  Acme is a
mouse-driven editor, and understanding the three buttons and their
chords is essential.

Chapter~\ref{ch:edit} (The Edit Command) covers the most powerful
feature---a complete structural editing language inherited from the sam
editor.  It rewards study but is not required for day-to-day use.

The appendices provide quick-reference tables for commands, mouse
actions, keyboard shortcuts, and the Unicode composition system.

Throughout this book, mouse buttons are referred to as \button{1}
(left), \button{2} (middle), and \button{3} (right).  Keyboard
shortcuts are written as \key{Ctrl+A} (hold Control, press A).
Built-in commands appear as \cmd{Put}, \cmd{Get}, \cmd{New}, etc.

\chapter{Installation}
\label{ch:install}

\index{installation}

Acme builds from source on Linux, macOS, and Windows.  The build
produces a single binary, \file{bin/acme}, with no runtime
configuration files required.

\section{Linux}
\label{sec:install-linux}

\index{installation!Linux}

\subsection{Prerequisites}

You need a C compiler (GCC or Clang), GNU make, and the development
libraries for SDL3, FreeType, fontconfig, and zlib.

\paragraph{Debian / Ubuntu:}
\begin{lstlisting}[language=bash]
sudo apt install gcc make
sudo apt install libsdl3-dev libfontconfig1-dev \
                 libfreetype-dev zlib1g-dev
\end{lstlisting}

\paragraph{Fedora:}
\begin{lstlisting}[language=bash]
sudo dnf install gcc make
sudo dnf install SDL3-devel fontconfig-devel \
                 freetype-devel zlib-devel
\end{lstlisting}

\paragraph{Arch Linux:}
\begin{lstlisting}[language=bash]
sudo pacman -S gcc make sdl3 fontconfig freetype2 zlib
\end{lstlisting}

Note: SDL3 is relatively new.  If your distribution does not yet
package it, you may need to build SDL3 from source or use a PPA.

\subsection{Building}

\begin{lstlisting}[language=bash]
git clone <repository-url> acme
cd acme
make
\end{lstlisting}

The build compiles 18 static libraries and the acme binary.  The result
is \file{bin/acme}.

\subsection{Running}

\begin{lstlisting}[language=bash]
./bin/acme
\end{lstlisting}

You can pass files to open on the command line:

\begin{lstlisting}[language=bash]
./bin/acme main.c util.c README.md
\end{lstlisting}

To install system-wide, copy the binary:

\begin{lstlisting}[language=bash]
sudo cp bin/acme /usr/local/bin/
\end{lstlisting}

\section{macOS}
\label{sec:install-macos}

\index{installation!macOS}

\subsection{Prerequisites}

Install Xcode Command Line Tools and Homebrew packages:

\begin{lstlisting}[language=bash]
xcode-select --install
brew install sdl3 fontconfig freetype pkg-config
\end{lstlisting}

\subsection{Building and Running}

\begin{lstlisting}[language=bash]
make
./bin/acme
\end{lstlisting}

The build uses \cmd{pkg-config} to locate Homebrew libraries.  If
\cmd{pkg-config} is not available, it falls back to default Homebrew
paths (\file{/opt/homebrew} on Apple Silicon, \file{/usr/local} on
Intel Macs).

\subsection{Mouse on macOS}

\index{mouse!macOS}
Most Mac users do not have a three-button mouse.  Acme provides
modifier-click emulation:

\begin{itemize}
\item \key{Ctrl}-click = \button{2} (middle button / execute)
\item \key{Alt/Option}-click = \button{3} (right button / look)
\end{itemize}

A standard USB or Bluetooth three-button mouse is strongly recommended
for extended use.  See Chapter~\ref{ch:mouse} for details.

\section{Windows}
\label{sec:install-windows}

\index{installation!Windows}

Acme on Windows can be built two ways:

\begin{itemize}
\item \textbf{MSYS2 build} — uses the MSYS2 POSIX compatibility layer.
  Easier to build but requires MSYS2 installed at runtime.
\item \textbf{Native standalone build} — produces a single
  \file{acme.exe} that runs on any Windows~10 or later machine with
  no other files required.
\end{itemize}

Both builds require MSYS2 at \emph{build} time for the toolchain.

\subsection{Prerequisites}

\begin{enumerate}
\item Install MSYS2 from \url{https://www.msys2.org/}.
\item Open any MSYS2 shell from the Start menu (MSYS, MINGW64, or
      UCRT64 — the build system pins the right compilers automatically).
\item Install dependencies:
\end{enumerate}

\begin{lstlisting}[language=bash]
pacman -Syu
pacman -S git gcc make pkg-config zlib-devel
pacman -S mingw-w64-x86_64-gcc \
          mingw-w64-x86_64-sdl3 \
          mingw-w64-x86_64-fontconfig \
          mingw-w64-x86_64-freetype
\end{lstlisting}

\subsection{Building (MSYS2 build)}

From any MSYS2 shell:

\begin{lstlisting}[language=bash]
export PKG_CONFIG_PATH=/mingw64/lib/pkgconfig
export PKG_CONFIG_ALLOW_SYSTEM_CFLAGS=1
export PKG_CONFIG_ALLOW_SYSTEM_LIBS=1
make
./bin/acme
\end{lstlisting}

The binary depends on the MSYS2 runtime (\file{msys-2.0.dll}) and
mingw64 DLLs.  It must be run from an MSYS2 environment or with
those DLLs in the PATH.

See \file{Windows-MSYS.md} in the source tree for full details.

\subsection{Building (native standalone)}

From any MSYS2 shell:

\begin{lstlisting}[language=bash]
make native-win
\end{lstlisting}

This produces \file{bin/acme.exe} (approximately 13~MB), statically
linked with no MSYS2 or Cygwin runtime dependency.  Copy
\file{acme.exe} to any Windows~10 or later machine and run it
directly --- no installation, no other files required.

The default font is Consolas, included with Windows.  Fonts are
discovered automatically from \file{C:\textbackslash Windows\textbackslash Fonts}.
The \cmd{Win} command opens an interactive \cmd{cmd.exe} session
inside an acme window using the Windows ConPTY API.

See \file{Windows-native.md} for full details.

\subsection{Mouse on Windows}

The same modifier-click emulation is available:

\begin{itemize}
\item \key{Ctrl}-click = \button{2} (middle button)
\item \key{Alt}-click = \button{3} (right button)
\end{itemize}

Any standard three-button USB mouse works natively.

\section{Command-Line Options}
\label{sec:cmdline}

\index{command-line options}

\begin{longtable}{@{}lp{3.5in}@{}}
\toprule
Flag & Description \\
\midrule
\endhead
\cmd{-a} & Enable auto-indent for all windows at startup. \\
\cmd{-c \textit{N}} & Start with \textit{N} columns (default: 2). \\
\cmd{-f \textit{font}} & Set the proportional (variable-width) font. \\
\cmd{-F \textit{font}} & Set the fixed-width font. \\
\cmd{-l \textit{file}} & Load session state from a dump file. \\
\cmd{-r} & Reverse scrollbar button mapping (Button~1 scrolls down). \\
\cmd{-W \textit{WxH}} & Set initial window size, e.g., \cmd{-W 1200x800}. \\
\bottomrule
\end{longtable}

\noindent
Fonts are specified as fontconfig names with a size suffix.  The
default for both \cmd{-f} and \cmd{-F} is
\file{DejaVuSansMono/13a/font}.

\section{Environment Variables}
\label{sec:envvars}

\index{environment variables}

\begin{longtable}{@{}lp{3.3in}@{}}
\toprule
Variable & Description \\
\midrule
\endhead
\env{HOME} & User home directory.  Used to find \file{\textasciitilde/lib/plumbing}. \\
\env{acmeshell} & Shell used to run commands.  Defaults to \cmd{sh}. \\
\env{tabstop} & Default tab width in columns.  Default is 4. \\
\env{SHELL} & Shell used by the Win command. \\
\env{NAMESPACE} & Plan~9 namespace directory.  Usually auto-detected. \\
\bottomrule
\end{longtable}

\noindent
When acme runs an external command, it sets several variables in the
child's environment:

\begin{longtable}{@{}lp{3.3in}@{}}
\toprule
Variable & Description \\
\midrule
\endhead
\env{winid} & Numeric ID of the current window. \\
\env{\%} & Filename of the current window. \\
\env{samfile} & Same as \env{\%}. \\
\bottomrule
\end{longtable}

\chapter{The Acme Screen}
\label{ch:screen}

\index{screen layout}

When acme starts, you see a deceptively simple display.  There are no
menus, no toolbars, no status bar, no file tree.  There is only text,
arranged in columns and windows.  Every element on screen is either
text you can edit or a border you can drag.

\section{Anatomy of the Display}

The acme screen is divided into a hierarchy:

\begin{enumerate}
\item A \textbf{top-level tag bar} spans the full width of the screen.
\item Below it, one or more \textbf{columns} are arranged side by side.
\item Each column contains one or more \textbf{windows}, stacked vertically.
\item Each window has its own \textbf{tag bar} and \textbf{body}.
\end{enumerate}

\subsection{The Top Tag Bar}
\index{tag bar!top-level}

The top tag bar shows global information and commands.  A typical top
tag bar reads:

\begin{Verbatim}[frame=single]
Newcol Kill Putall Dump Exit
\end{Verbatim}

These are not buttons---they are editable text.  You can add your own
commands by typing them into the tag bar.  Middle-click (\button{2}) on
any word to execute it.

\subsection{Columns}
\index{columns}

Columns are vertical panes that divide the screen horizontally.  Each
column has a thin tag bar at its top, typically containing:

\begin{Verbatim}[frame=single]
New Cut Paste Snarf Sort Zerox Delcol
\end{Verbatim}

\noindent
You can:

\begin{itemize}
\item Create a new column by executing \cmd{Newcol} in the top tag bar.
\item Delete a column by executing \cmd{Delcol} in the column's tag bar.
\item Resize columns by dragging the thin border between them with
      \button{1}.
\end{itemize}

\subsection{Windows}
\index{windows}

Each window has two parts:

\paragraph{The tag bar}
\index{tag bar!window}
is a single line (or more, if the content wraps) at the top of the
window.  It contains the window's filename followed by commands.  A
typical tag bar reads:

\begin{Verbatim}[frame=single]
/home/user/src/main.c Del Snarf Get Put Undo Redo | Look
\end{Verbatim}

The filename at the left is editable---you can change it to rename the
file, or clear it for a scratch buffer.  Everything after the vertical
bar \texttt{|} is the ``user area'' of the tag, where you can type
whatever commands you find useful.

\paragraph{The body}
\index{body}
occupies the rest of the window and contains the file's text.  This is
where you read and edit.

\subsection{The Layout Button}
\index{layout button}

At the far left of each window's tag bar is a small square.  This is
the \emph{layout button}.  It serves two purposes:

\begin{itemize}
\item Drag it with \button{1} to move the window within its column or
      to another column.
\item Its appearance indicates the window's state:
  \begin{itemize}
  \item A clean square means the window matches the file on disk.
  \item A filled (dirty) square means the window has unsaved changes.
  \end{itemize}
\end{itemize}

\subsection{The Scrollbar}
\index{scrollbar}

Along the left edge of each window body is a narrow scrollbar.  It
shows a small rectangle (the ``thumb'') indicating which portion of the
file is currently visible.  See Section~\ref{sec:scrollbar} for details
on mouse interaction with the scrollbar.

\section{Text Everywhere}

The key insight of acme's interface is that \emph{everything is text}.
The tag bars are text.  The body is text.  Commands are text.  You
interact with all of it the same way: by clicking on it with the mouse.

This means:

\begin{itemize}
\item You can type a command anywhere---in a tag, in a body, even in
      the middle of a file---and middle-click it to execute.
\item You can type a filename in any window and right-click it to open
      that file.
\item Output from commands appears as text in windows, which you can
      then edit, search, or use as input to other commands.
\end{itemize}

There is no distinction between ``code,'' ``output,'' and ``interface.''
It is all text, and it is all interactive.

\section{Window Management}

\index{window management}

Windows in acme are not managed by tabs or a file tree.  Instead:

\begin{itemize}
\item Open a file by right-clicking (\button{3}) its name.
\item Create an empty window with the \cmd{New} command.
\item Close a window with the \cmd{Del} command.
\item Clone a window with the \cmd{Zerox} command.
\item Sort windows alphabetically with the \cmd{Sort} command.
\item Move windows by dragging the layout button.
\end{itemize}

To navigate between windows, you simply click in the one you want.
There is no ``active window'' highlight or focus indicator---you click
where you want to work, and that window responds.

\section{Directory Windows}
\index{directory windows}

Right-clicking a directory path opens a window listing the directory's
contents.  Each entry in the listing is itself clickable: right-click a
filename to open it, right-click a subdirectory to list it.

Directory windows update automatically when files change.  They serve
the role of a file browser, but they are just text in a window---you
can search them, select from them, or pipe their contents through
commands.

\section{The +Errors Window}
\index{+Errors window}
\index{errors}

When an external command writes to standard error, the output appears in
a window named \file{+Errors} (prefixed with the directory path).
These windows are created automatically as needed.  You can close them
with \cmd{Del} like any other window.

Compiler errors often appear in the format \file{file.c:27: error...}.
Right-clicking (\button{3}) on such a line opens \file{file.c} at
line~27---a simple but powerful way to navigate error output.

\chapter{The Mouse}
\label{ch:mouse}

\index{mouse}

This is the most important chapter in this book.  Acme is a
mouse-driven editor.  The three mouse buttons are the primary way you
interact with acme---more so than the keyboard.  If you learn nothing
else, learn the three buttons and the chords.

Throughout this chapter, the buttons are:

\begin{description}
\item[\button{1}] Left button --- \textbf{Select}
\item[\button{2}] Middle button --- \textbf{Execute}
\item[\button{3}] Right button --- \textbf{Look / Open}
\end{description}

\section{Button 1: Select}
\label{sec:b1}

\index{mouse!Button 1}
\index{selection}

\button{1} is for selecting text.

\begin{description}[style=nextline]
\item[Click]
Places the cursor (an empty selection) at the click position.

\item[Drag]
Selects a range of text.  The selected text is highlighted.

\item[Double-click]
Selects the word under the cursor.  A ``word'' is a run of
alphanumeric characters and underscores, or a run of
non-whitespace characters, depending on context.

\item[Double-click on a bracket]
If you double-click on any of \texttt{( ) [ ] \{ \} < >}, acme
selects the text enclosed by the matching pair, including nested
brackets.  This works for parentheses, square brackets, curly
braces, and angle brackets.

\item[Double-click on a quote]
Double-clicking on a single quote, double quote, or backtick
selects the text enclosed by the matching quote character.
\end{description}

Clicking \button{1} anywhere deselects the previous selection.  The
selection is fundamental to acme: many commands operate on the selected
text, and the selected text is what gets cut, copied, or replaced.

\section{Button 2: Execute}
\label{sec:b2}

\index{mouse!Button 2}
\index{executing commands}

\button{2} executes text as a command.

\begin{description}[style=nextline]
\item[Click on a built-in command]
Middle-click on \cmd{Put} and acme saves the file.  Middle-click on
\cmd{Del} and acme closes the window.  Middle-click on \cmd{Undo}
and acme undoes the last change.

\item[Click on arbitrary text]
Middle-click on any text to run it as a shell command.  For example,
type \cmd{date} in a tag bar, middle-click it, and the current date
appears in the \file{+Errors} window.

\item[Click on a blank area]
If you click in an empty area or on whitespace, acme expands the
selection to the surrounding non-whitespace text and executes that.
This means you do not need to carefully select a command word---just
click anywhere on it.

\item[Click on a pipe command]
Text beginning with \texttt{<}, \texttt{>}, or \texttt{|} has special
meaning.  See Chapter~\ref{ch:shell} for details.
\end{description}

The command executes in the directory associated with the window.  For a
file \file{/home/user/src/main.c}, commands run in
\file{/home/user/src/}.

\section{Button 3: Look and Open}
\label{sec:b3}

\index{mouse!Button 3}
\index{Look}
\index{plumbing}

\button{3} is the ``smart open'' button.  It examines the text you
click on and does the most useful thing:

\begin{description}[style=nextline]
\item[Filename]
Right-click on \file{main.c} and acme opens it in a new window (or
jumps to the existing window if it is already open).

\item[Filename with address]
Right-click on \file{main.c:42} and acme opens \file{main.c} and
jumps to line~42.  The address can be a line number, a line:column
pair (\file{main.c:42.10}), or a regular expression
(\file{main.c:/pattern/}).

\item[Address in current file]
Right-click on \texttt{:42} or \texttt{:/pattern/} to jump to that
location in the current window.

\item[Include file]
Right-click on \texttt{<stdio.h>} and acme searches include
directories for \file{stdio.h} and opens it.

\item[URL]
Right-click on \texttt{https://example.com} and acme opens it in
your web browser (via \cmd{xdg-open}, \cmd{open}, or
\cmd{cygstart} depending on platform).

\item[Plain text]
Right-click on any other text and acme searches forward in the
current window body for that text---a literal search (not a regex).
Click again to find the next occurrence.
\end{description}

The decision-making process is handled by the \emph{plumber} (see
Chapter~\ref{ch:plumbing}).  The built-in plumbing rules handle
filenames, URLs, images, and PDFs.  You can add custom rules.

\section{Mouse Chords}
\label{sec:chords}

\index{mouse!chords}
\index{chords}

Mouse chords are combinations of buttons pressed in sequence.  They
provide quick access to cut, paste, and copy without reaching for menu
commands.

\subsection{Cut: \chord{1}{2}}
\index{Cut!chord}

\begin{enumerate}
\item Press and hold \button{1}, drag to select text.
\item While still holding \button{1}, click \button{2}.
\item Release both buttons.
\end{enumerate}

The selected text is deleted and placed in the snarf buffer (clipboard).
This is identical to executing the \cmd{Cut} command.

\subsection{Paste: \chord{1}{3}}
\index{Paste!chord}

\begin{enumerate}
\item Press and hold \button{1} to position the cursor or select text.
\item While still holding \button{1}, click \button{3}.
\item Release both buttons.
\end{enumerate}

The snarf buffer contents replace the selection (or are inserted at the
cursor position).  This is identical to executing the \cmd{Paste}
command.

\subsection{Snarf: \chord{1}{2} then \button{3}}
\index{Snarf!chord}

\begin{enumerate}
\item Press and hold \button{1}, drag to select text.
\item While holding \button{1}, click \button{2} (this cuts).
\item While still holding \button{1}, click \button{3} (this pastes
      back what was just cut).
\item Release all buttons.
\end{enumerate}

The net effect: the text remains unchanged, but it has been copied to
the snarf buffer.  This is the chord for ``copy without deleting.''

\subsection{The 2-1 Argument Chord}
\index{2-1 chord}

Commands can receive an argument via a two-button chord:

\begin{enumerate}
\item Press and hold \button{2} on a command word (e.g., \cmd{Look}).
\item While holding \button{2}, click \button{1} on the argument text.
\item Release both buttons.
\end{enumerate}

Acme passes the \button{1}-selected text as an argument to the
\button{2} command.  For example, holding \button{2} on \cmd{Look}
and clicking \button{1} on a word searches for that word.

\section{The Scrollbar}
\label{sec:scrollbar}

\index{scrollbar}

The scrollbar is the narrow strip along the left edge of each window
body.  It contains a thumb (small rectangle) showing the current
viewport position relative to the whole file.

\begin{description}[style=nextline]
\item[\button{1} in scrollbar]
Scrolls \textbf{up}.  The line at the click position moves to the
top of the window.

\item[\button{2} in scrollbar]
\textbf{Jump scroll}.  The display jumps so that the thumb is at the
click position---proportional scrolling.

\item[\button{3} in scrollbar]
Scrolls \textbf{down}.  The line at the top of the window moves to
the click position.
\end{description}

The \cmd{-r} command-line flag reverses the sense of \button{1} and
\button{3} in the scrollbar, for users who prefer the opposite
convention.

\section{Scroll Wheel}
\label{sec:scrollwheel}

\index{scroll wheel}

The scroll wheel scrolls the window body:

\begin{itemize}
\item Scroll up = scroll text up (view earlier content).
\item Scroll down = scroll text down (view later content).
\end{itemize}

The scroll amount is a fraction of the visible window height.

\section{One-Button and Two-Button Mice}
\label{sec:emulation}

\index{mouse!emulation}
\index{Ctrl-click}
\index{Alt-click}

On laptops, trackpads, and Mac mice, you may not have three distinct
buttons.  Acme provides modifier-click emulation:

\begin{longtable}{@{}ll@{}}
\toprule
Modifier Click & Equivalent \\
\midrule
\endhead
\key{Ctrl}-click & \button{2} (execute) \\
\key{Alt}-click (or \key{Option}-click on Mac) & \button{3} (look/open) \\
\bottomrule
\end{longtable}

These modifiers work with \button{1} (left click).  So on a trackpad:

\begin{itemize}
\item Plain click = \button{1} (select)
\item \key{Ctrl}+click = \button{2} (execute)
\item \key{Alt}+click = \button{3} (look/open)
\end{itemize}

\textbf{Recommendation:} For serious use, get a three-button mouse.
Acme was designed for one, and the experience is dramatically better
with real buttons.  Any inexpensive USB mouse with a clickable scroll
wheel provides three buttons.

\section{Summary}

\begin{longtable}{@{}llp{2.5in}@{}}
\toprule
Button & Name & Action \\
\midrule
\endhead
\button{1} & Select & Select text, position cursor, drag to move windows \\
\button{2} & Execute & Run built-in commands or shell commands \\
\button{3} & Look & Open files, search text, follow addresses \\
\chord{1}{2} & Cut & Delete selection to snarf buffer \\
\chord{1}{3} & Paste & Replace selection with snarf buffer \\
\chord{1}{2}+\button{3} & Snarf & Copy selection to snarf buffer \\
\bottomrule
\end{longtable}

\chapter{Keyboard}
\label{ch:keyboard}

\index{keyboard}

While the mouse is acme's primary interaction tool, the keyboard
handles text entry, cursor movement, and a small set of editing
shortcuts.  This standalone version of acme remaps the cursor keys,
Home, End, and Page keys to work in a standard, familiar way.

\section{Cursor Movement}
\label{sec:cursor}

\index{cursor movement}

\begin{longtable}{@{}lp{3.2in}@{}}
\toprule
Key & Action \\
\midrule
\endhead
\key{Left Arrow} & Move cursor one character left. \\
\key{Right Arrow} & Move cursor one character right. \\
\key{Up Arrow} & Move cursor up one line. \\
\key{Down Arrow} & Move cursor down one line. \\
\key{Ctrl+Left} & Move to beginning of line (same as \key{Ctrl+A}). \\
\key{Ctrl+Right} & Move to end of line (same as \key{Ctrl+E}). \\
\key{Ctrl+A} & Move to beginning of line. \\
\key{Ctrl+E} & Move to end of line. \\
\bottomrule
\end{longtable}

\subsection{Sticky Column}
\index{sticky column}

The up and down arrow keys maintain a \emph{sticky column} position.
When you press \key{Up} or \key{Down}, acme remembers the horizontal
pixel position of the cursor and tries to place it at the same
horizontal position on the adjacent line.  If the adjacent line is
shorter, the cursor goes to the end of that line, but the original
column is remembered.  When you reach a longer line again, the cursor
returns to the remembered column.

The sticky column resets whenever you press any key other than \key{Up}
or \key{Down}, or when you click the mouse.

\subsection{Scrolling with Arrow Keys}
\index{scrolling!arrow keys}

If the cursor is on the \emph{last visible line} and you press
\key{Down}, acme scrolls down one line and places the cursor on the new
last visible line.  Similarly, pressing \key{Up} on the first visible
line scrolls up one line.  This gives smooth, continuous scrolling while
navigating with arrow keys.

\section{Page Navigation}
\label{sec:paging}

\index{scrolling!page}

\begin{longtable}{@{}lp{3.2in}@{}}
\toprule
Key & Action \\
\midrule
\endhead
\key{Page Up} & Scroll up one full screen. \\
\key{Page Down} & Scroll down one full screen. \\
\key{Ctrl+Up} & Scroll up one full screen (same as Page~Up). \\
\key{Ctrl+Down} & Scroll down one full screen (same as Page~Down). \\
\bottomrule
\end{longtable}

\key{Page Down} does nothing if the end of the file is already visible.

\section{Jump to Position}
\label{sec:jump}

\index{Home key}
\index{End key}

\begin{longtable}{@{}lp{3.2in}@{}}
\toprule
Key & Action \\
\midrule
\endhead
\key{Home} & Move cursor to first character visible on screen. \\
\key{End} & Move cursor to beginning of last visible line. \\
\key{Ctrl+Home} & Jump to beginning of file. \\
\key{Ctrl+End} & Jump to end of file. \\
\key{Ctrl+Page Up} & Jump to beginning of file (same as \key{Ctrl+Home}). \\
\key{Ctrl+Page Down} & Jump to end of file (same as \key{Ctrl+End}). \\
\bottomrule
\end{longtable}

\key{Home} and \key{End} do \emph{not} scroll the text---they move the
cursor within the currently visible portion.  \key{Ctrl+Home},
\key{Ctrl+End}, \key{Ctrl+Page Up}, and \key{Ctrl+Page Down} scroll
as needed to show the target position.

\section{Editing Keys}
\label{sec:editkeys}

\index{Backspace}
\index{Delete key}

\begin{longtable}{@{}lp{3.2in}@{}}
\toprule
Key & Action \\
\midrule
\endhead
\key{Backspace} & Delete the character to the left of the cursor. \\
\key{Delete} & Delete the character to the right of the cursor. \\
\key{Ctrl+H} & Delete character left (same as \key{Backspace}). \\
\key{Ctrl+W} & Delete word to the left of the cursor. \\
\key{Ctrl+U} & Delete from cursor to beginning of line. \\
\bottomrule
\end{longtable}

\key{Backspace} deletes the character to the left, which is the
standard behavior on modern systems.  \key{Delete} (forward delete)
removes the character to the right.

\section{Text Entry}

\index{auto-indent}
\index{Tab key}

Typing inserts text at the cursor position, replacing any selection.
This is standard behavior---there are no modes.

\paragraph{Auto-indent.}
When auto-indent is enabled (via the \cmd{-a} flag or the \cmd{Indent}
command), pressing \key{Enter} inserts a newline followed by the same
leading whitespace as the previous line.

\paragraph{Tab.}
The \key{Tab} key inserts a literal tab character.  The display width of
a tab is controlled by the \cmd{Tab} command or the \env{tabstop}
environment variable (default: 4 columns).

\section{Unicode Composition}
\label{sec:compose}

\index{Unicode composition}
\index{Alt key}

Acme supports entering Unicode characters through a multi-key
composition system.  To compose a character:

\begin{enumerate}
\item Press and release the \key{Alt} key (this toggles composition mode).
\item Type the composition sequence (one or more characters).
\item The composed character is inserted.
\end{enumerate}

If the sequence is unrecognized, the literal characters are inserted.
Press \key{Alt} again to cancel composition mode.

\subsection{Common Compositions}

\begin{longtable}{@{}llll@{}}
\toprule
Sequence & Result & & Description \\
\midrule
\endhead
\texttt{'e} & \'{e} & & Acute accent \\
\texttt{`e} & \`{e} & & Grave accent \\
\texttt{\^{}e} & \^{e} & & Circumflex \\
\texttt{"e} & \"{e} & & Diaeresis / umlaut \\
\texttt{\~{}n} & \~{n} & & Tilde \\
\texttt{,c} & \c{c} & & Cedilla \\
\texttt{oe} & \oe & & Ligature \\
\texttt{*a} & $\alpha$ & & Greek alpha \\
\texttt{*b} & $\beta$ & & Greek beta \\
\texttt{*p} & $\pi$ & & Greek pi \\
\texttt{!=} & $\neq$ & & Not equal \\
\texttt{<=} & $\leq$ & & Less or equal \\
\texttt{>=} & $\geq$ & & Greater or equal \\
\texttt{+-} & $\pm$ & & Plus-minus \\
\texttt{->} & $\rightarrow$ & & Right arrow \\
\texttt{<-} & $\leftarrow$ & & Left arrow \\
\texttt{sr} & $\sqrt{}$ & & Square root \\
\texttt{if} & $\infty$ & & Infinity \\
\bottomrule
\end{longtable}

\noindent
The full composition table contains over 500 entries, covering Latin
accented characters, Greek, Cyrillic, mathematical symbols, currency
symbols, arrows, and more.  See Appendix~\ref{app:unicode} for the
complete table.

\section{Summary}

\begin{longtable}{@{}llp{2.5in}@{}}
\toprule
Category & Keys & Notes \\
\midrule
\endhead
Movement & Arrows, \key{Ctrl}+Arrows & Sticky column on up/down \\
Pages & \key{PgUp}, \key{PgDn}, \key{Ctrl}+Up/Dn & Full-screen scroll \\
Jump & \key{Home}, \key{End} & Within visible area \\
File jump & \key{Ctrl+Home}, \key{Ctrl+End} & Beginning / end of file \\
Delete & \key{Backspace}, \key{Delete} & Left / right of cursor \\
Words & \key{Ctrl+W} & Delete word left \\
Lines & \key{Ctrl+U}, \key{Ctrl+A}, \key{Ctrl+E} & Delete/move to line boundary \\
Compose & \key{Alt}, then sequence & Unicode character entry \\
\bottomrule
\end{longtable}

\chapter{Built-in Commands}
\label{ch:commands}

\index{commands}

Acme's built-in commands are capitalized words that you execute by
middle-clicking (\button{2}) on them.  They typically appear in window
tag bars but can be typed and executed anywhere.

\section{Clipboard: Cut, Snarf, Paste}
\label{sec:clipboard}

\index{Cut}
\index{Snarf}
\index{Paste}
\index{snarf buffer}
\index{clipboard}

Acme maintains an internal buffer called the \emph{snarf buffer}.  This
is also synchronized with the system clipboard via SDL3, so you can
copy and paste between acme and other applications.

\begin{description}[style=nextline]
\item[\cmd{Cut}]
Delete the selected text and place it in the snarf buffer.  The
window is marked dirty.

\item[\cmd{Snarf}]
Copy the selected text to the snarf buffer \emph{without} deleting
it.  The window is not marked dirty.

\item[\cmd{Paste}]
Replace the current selection with the contents of the snarf buffer.
If there is no selection, insert at the cursor position.
\end{description}

The mouse chords \chord{1}{2} (Cut), \chord{1}{3} (Paste), and
\chord{1}{2}+\button{3} (Snarf) perform the same operations more
quickly.  See Section~\ref{sec:chords}.

To paste text from another application into acme: copy in the other
application (e.g., \key{Ctrl+C}), then execute \cmd{Paste} in acme or
use the \chord{1}{3} chord.

\section{Undo and Redo}
\label{sec:undo}

\index{Undo}
\index{Redo}

\begin{description}[style=nextline]
\item[\cmd{Undo}]
Undo the last change or group of changes.  Acme groups related
changes (e.g., all characters typed before a pause) into a single
undo step.

\item[\cmd{Redo}]
Reverse the last \cmd{Undo}.
\end{description}

Undo operates on a per-file basis.  If the same file is open in
multiple windows, undoing in one window affects all windows showing that
file.

\section{File Operations: Get, Put, Putall}
\label{sec:fileops}

\index{Get}
\index{Put}
\index{Putall}

\begin{description}[style=nextline]
\item[\cmd{Get}]
Reload the file from disk, replacing the window's contents.  If the
window has unsaved changes, acme warns you (execute \cmd{Get} a
second time to override).

With an argument (\cmd{Get \textit{filename}}), load the named file
into the current window without changing the window's name.

\item[\cmd{Put}]
Write the window's contents to disk.  Uses the filename shown in the
tag bar.

With an argument (\cmd{Put \textit{filename}}), write to the named
file instead (``Save As'').

\item[\cmd{Putall}]
Write all dirty windows to their respective files.  Only saves
windows whose names correspond to real files.  Skips scratch
windows, directory windows, and windows with no filename.
\end{description}

\section{Window Management}
\label{sec:winmgmt}

\index{New}
\index{Del}
\index{Delete}
\index{Zerox}
\index{Newcol}
\index{Delcol}
\index{Sort}

\begin{description}[style=nextline]
\item[\cmd{New}]
Create a new empty window.  With arguments (\cmd{New
\textit{file1 file2...}}), open those files in new windows.

\item[\cmd{Del}]
Close the current window.  If the window has unsaved changes, acme
warns you.  Execute \cmd{Del} a second time to close anyway.

\item[\cmd{Delete}]
Close the current window without checking for unsaved changes.

\item[\cmd{Zerox}]
Clone the current window, creating a second view of the same file.
Both windows show the same underlying file; changes in one appear in
the other.

\item[\cmd{Newcol}]
Create a new column.

\item[\cmd{Delcol}]
Delete the current column and all its windows.  Warns if any window
has unsaved changes.

\item[\cmd{Sort}]
Sort the windows in the current column alphabetically by filename.
\end{description}

\section{Session Management: Dump, Load, Exit}
\label{sec:session}

\index{Dump}
\index{Load}
\index{Exit}

\begin{description}[style=nextline]
\item[\cmd{Dump}]
Save the complete acme session state---window layout, filenames, font
settings, and positions---to a file.  Default: \file{\$HOME/acme.dump}.

With an argument (\cmd{Dump \textit{filename}}), save to the named
file.

\item[\cmd{Load}]
Restore a previously saved session from a dump file.  Default:
\file{\$HOME/acme.dump}.

\item[\cmd{Exit}]
Quit acme.  If any window has unsaved changes, acme warns you.
Execute \cmd{Exit} a second time to quit anyway.
\end{description}

The \cmd{Dump}/\cmd{Load} pair lets you preserve your workspace
across sessions.  It is common to dump before quitting and load when
restarting.

\section{Search: Look}
\label{sec:look}

\index{Look}

\begin{description}[style=nextline]
\item[\cmd{Look}]
Search for text in the window body.  Uses the current selection as
the search term.  If there is no selection, uses the word nearest
the cursor.

With an argument (via the 2-1 chord or by typing \cmd{Look
\textit{text}}), search for that text.

The search is \textbf{literal} (not a regex) and wraps around the
end of the file.  For regex search, use the \cmd{Edit} command
(Chapter~\ref{ch:edit}) or right-click on a \texttt{:/regex/}
address.
\end{description}

\section{Display: Font, Tab, Indent}
\label{sec:display}

\index{Font}
\index{Tab}
\index{Indent}
\index{auto-indent}

\begin{description}[style=nextline]
\item[\cmd{Font}]
Toggle between proportional and fixed-width fonts.

With an argument (\cmd{Font \textit{fontname}}), set the window's
font.  Font names are fontconfig names with a size suffix (e.g.,
\file{DejaVuSans/12a/font}).

\item[\cmd{Tab}]
With no argument: print the current tab width.

With a numeric argument (\cmd{Tab 8}): set the tab display width for
the current window.

\item[\cmd{Indent}]
Control auto-indentation.

\begin{itemize}
\item \cmd{Indent on} / \cmd{Indent off} --- affect the current window.
\item \cmd{Indent ON} / \cmd{Indent OFF} --- affect all windows globally.
\end{itemize}
\end{description}

\section{Process Control: Kill, ID}
\label{sec:process}

\index{Kill}
\index{ID}

\begin{description}[style=nextline]
\item[\cmd{Kill}]
Send a kill signal to running commands.  With arguments, kills
commands whose names match.

\item[\cmd{ID}]
Print the current window's numeric ID.
\end{description}

\section{Include Paths: Incl}
\label{sec:incl}

\index{Incl}

\begin{description}[style=nextline]
\item[\cmd{Incl}]
With no arguments: list the current supplementary include directories.

With arguments (\cmd{Incl \textit{/path/to/headers}}): add
directories to the search path used when right-clicking on
\texttt{<filename.h>} patterns.
\end{description}

The default search path includes \file{/usr/include} and
\file{/usr/local/include}.  The \cmd{Incl} command adds project-specific
directories.

\section{Send}
\label{sec:send}

\index{Send}

\begin{description}[style=nextline]
\item[\cmd{Send}]
In a Win window (see Chapter~\ref{ch:shell}), send the text between
the prompt boundary and the end of the body to the shell.

In a regular window, append the snarf buffer contents to the end
of the body.
\end{description}

\section{Win}
\label{sec:win-brief}

\index{Win}

The \cmd{Win} command opens an interactive shell session in an acme
window.  See Chapter~\ref{ch:shell} for full details.

\chapter{Shell Commands and Pipes}
\label{ch:shell}

\index{shell commands}
\index{pipe commands}

One of acme's most powerful features is its seamless integration with
external commands.  Any text can be executed as a shell command.  The
Unix toolkit becomes your extension mechanism.

\section{Running External Commands}
\label{sec:external}

\index{executing commands!external}

When you middle-click (\button{2}) on text that is not a built-in
command, acme runs it as a shell command.  The shell used is determined
by the \env{acmeshell} environment variable, defaulting to \cmd{sh}.

The command runs in the directory associated with the current window.
For a file \file{/home/user/project/main.c}, commands execute in
\file{/home/user/project/}.

\subsection{Environment}

Acme sets the following environment variables for child processes:

\begin{itemize}
\item \env{winid} --- the numeric ID of the window.
\item \env{\%} --- the filename of the window.
\item \env{samfile} --- same as \env{\%}.
\end{itemize}

\subsection{Output}

Standard output and standard error from the command appear in the
\file{+Errors} window for that directory.  If no \file{+Errors} window
exists, one is created automatically.

\subsection{Examples}

Type any of these in a tag bar and middle-click to execute:

\begin{lstlisting}[language=bash]
date                     # show the date
ls -la                   # list files
make                     # build the project
grep -rn TODO *.c        # search for TODOs
git status               # check version control
\end{lstlisting}

\section{Pipe Prefixes: <, >, |}
\label{sec:pipes}

\index{pipe!< prefix}
\index{pipe!> prefix}
\index{pipe!| prefix}

Commands prefixed with \texttt{<}, \texttt{>}, or \texttt{|} interact
with the current selection:

\begin{longtable}{@{}cp{3.5in}@{}}
\toprule
Prefix & Behavior \\
\midrule
\endhead
\texttt{<cmd} & Run \texttt{cmd}, \textbf{replace the selection}
  with its standard output. \\
\texttt{>cmd} & Run \texttt{cmd}, sending the \textbf{selection as
  standard input}.  The selection is unchanged. \\
\texttt{|cmd} & Run \texttt{cmd}, sending the \textbf{selection as
  standard input} and \textbf{replacing it} with standard output. \\
\bottomrule
\end{longtable}

\subsection{Examples}

\paragraph{Replace selection with command output (\texttt{<}):}

\begin{lstlisting}[language=bash]
<date              # insert the current date
<seq 1 10          # insert numbers 1 through 10
<cat /etc/hostname # insert the hostname
\end{lstlisting}

\paragraph{Send selection to command (\texttt{>}):}

\begin{lstlisting}[language=bash]
>wc            # count words/lines in selection
>lp            # print selection
>xclip         # copy to X clipboard
>mail user@x   # email the selection
\end{lstlisting}

\paragraph{Filter selection through command (\texttt{\textbar}):}

\begin{lstlisting}[language=bash]
|sort             # sort selected lines
|sort -u          # sort and remove duplicates
|fmt -w 72        # reflow text to 72 columns
|tr a-z A-Z       # uppercase the selection
|awk '{print NR, $0}'  # add line numbers
|sed 's/old/new/g'     # search and replace
|column -t        # align columns
|indent           # reformat C code
|gofmt            # reformat Go code
\end{lstlisting}

These pipe commands are the primary way to extend acme.  Instead of
plugins, you use the standard Unix toolkit.

\section{The +Errors Window}
\label{sec:errors}

\index{+Errors window}

Command output goes to a window named with the directory path plus
\file{+Errors}.  For example, commands run from
\file{/home/user/project/main.c} produce output in
\file{/home/user/project/+Errors}.

The \file{+Errors} window is an ordinary acme window.  You can:

\begin{itemize}
\item Search it with \button{3}.
\item Right-click on error messages like \file{main.c:42: error} to
      jump directly to the source location.
\item Close it with \cmd{Del}.
\item Clear it by selecting all and deleting.
\end{itemize}

\section{The Win Command}
\label{sec:win}

\index{Win}
\index{PTY}
\index{interactive shell}

The \cmd{Win} command opens an interactive shell session inside an acme
window.  It creates a pseudo-terminal (PTY) connected to your shell
(typically bash or whatever \env{SHELL} is set to).

\subsection{How It Works}

\begin{enumerate}
\item Execute \cmd{Win} (middle-click it in any tag bar).
\item A new window opens, named after the current directory with a
      \file{-win} suffix.
\item Type commands and press \key{Enter} to send them to the shell.
\item Output from the shell appears in the window body.
\end{enumerate}

\subsection{Behavior}

The Win window is not a terminal emulator---it is an acme text buffer
connected to a shell via PTY.  This means:

\begin{itemize}
\item You can edit text before pressing \key{Enter}.  The shell does not
      see your keystrokes until you press \key{Enter}.
\item You can select output text and use it like any other acme text
      (right-click filenames, middle-click commands, etc.).
\item ANSI escape sequences are stripped.  Colors and cursor positioning
      do not work.
\item The PTY is set to single-column width, so \cmd{ls} and similar
      commands produce one-entry-per-line output.
\item Echo is disabled---acme shows what you type, not the terminal.
\end{itemize}

\subsection{Shell Configuration}

The Win command starts the shell with a minimal environment:

\begin{itemize}
\item \texttt{TERM=dumb} --- suppresses escape sequences.
\item \texttt{PS1="\$ "} --- a simple prompt.
\item \texttt{PROMPT\_COMMAND}, \texttt{COLORTERM}, and
      \texttt{LS\_COLORS} are unset.
\end{itemize}

The shell is started with \cmd{--noediting --norc --noprofile -i}
flags (for bash) to minimize interference.

\subsection{Closing}

Close a Win window with \cmd{Del} like any other window.  The shell
process receives SIGHUP and terminates.

\section{When to Use What}

\begin{longtable}{@{}lp{3.5in}@{}}
\toprule
Need & Approach \\
\midrule
\endhead
Run a one-off command & Type it in the tag, \button{2} click \\
Transform selected text & Use \texttt{|cmd} \\
Insert command output & Use \texttt{<cmd} \\
Send text to a command & Use \texttt{>cmd} \\
Interactive session & Use \cmd{Win} \\
Navigate compiler errors & Right-click error lines in \file{+Errors} \\
\bottomrule
\end{longtable}

\chapter{The Edit Command}
\label{ch:edit}

\index{Edit command}
\index{sam}
\index{structural regular expressions}

The \cmd{Edit} command gives acme a complete text-processing language,
inherited from Rob Pike's sam editor.  It is the most powerful feature
in acme and rewards study, though it is not required for daily use.

To use it, type \cmd{Edit} followed by a sam command in any tag or
body, and middle-click (\button{2}) the whole thing.  For example:

\begin{Verbatim}[frame=single]
Edit ,s/old/new/g
\end{Verbatim}

\noindent
This replaces all occurrences of ``old'' with ``new'' in the current
window.

\section{Addresses}
\label{sec:addresses}

\index{addresses}

An address identifies a substring of the file.  Addresses appear before
commands to specify what text the command operates on.

\begin{longtable}{@{}lp{3.5in}@{}}
\toprule
Address & Meaning \\
\midrule
\endhead
\texttt{.} & Dot: the current selection. \\
\texttt{0} & The empty string at the beginning of the file. \\
\texttt{\$} & The empty string at the end of the file. \\
\texttt{\textit{n}} & Line \textit{n}. \\
\texttt{\#\textit{n}} & Character position \textit{n} (counting from 0). \\
\texttt{/\textit{regexp}/} & Next match of \textit{regexp} after dot. \\
\texttt{?\textit{regexp}?} & Previous match of \textit{regexp} before dot. \\
\texttt{\textit{a},\textit{b}} & From address \textit{a} to address \textit{b}. \\
\texttt{\textit{a};\textit{b}} & From \textit{a} to \textit{b}, with dot set to \textit{a}
  before evaluating \textit{b}. \\
\texttt{\textit{a}+\textit{b}} & Address \textit{b} evaluated from the end of \textit{a}. \\
\texttt{\textit{a}-\textit{b}} & Address \textit{b} evaluated from the start of \textit{a},
  searching backward. \\
\bottomrule
\end{longtable}

\paragraph{Examples:}

\begin{itemize}
\item \texttt{1,5} --- lines 1 through 5.
\item \texttt{,} --- the entire file (shorthand for \texttt{0,\$}).
\item \texttt{.+1} --- the line after the current selection.
\item \texttt{/func/} --- the next occurrence of ``func''.
\item \texttt{0,/\^{}package/} --- from start of file to the line
      matching ``package''.
\end{itemize}

\section{Simple Commands}
\label{sec:simple}

\begin{longtable}{@{}lp{3.5in}@{}}
\toprule
Command & Action \\
\midrule
\endhead
\texttt{a/\textit{text}/} & \textbf{Append}: insert \textit{text}
  after the addressed text. \\
\texttt{i/\textit{text}/} & \textbf{Insert}: insert \textit{text}
  before the addressed text. \\
\texttt{c/\textit{text}/} & \textbf{Change}: replace the addressed
  text with \textit{text}. \\
\texttt{d} & \textbf{Delete}: delete the addressed text. \\
\texttt{s/\textit{regexp}/\textit{text}/} & \textbf{Substitute}:
  replace the first match of \textit{regexp} in the addressed text
  with \textit{text}. \\
\texttt{s/\textit{regexp}/\textit{text}/g} & Substitute all matches
  (global flag). \\
\texttt{m \textit{addr}} & \textbf{Move}: move the addressed text to
  after \textit{addr}. \\
\texttt{t \textit{addr}} & \textbf{Copy}: copy the addressed text to
  after \textit{addr}. \\
\texttt{p} & \textbf{Print}: display the addressed text (in
  \file{+Errors}). \\
\texttt{=} & Print the line number and character offset of the
  address. \\
\bottomrule
\end{longtable}

The delimiter for \texttt{a}, \texttt{i}, \texttt{c}, and \texttt{s}
can be any character, not just \texttt{/}.  Newlines within the text are
entered literally (the text continues on the next line, terminated by
another delimiter).

\paragraph{Examples:}

\begin{Verbatim}[frame=single]
Edit 1 i/// Copyright 2024\n/
Edit ,d
Edit ,s/colour/color/g
Edit 5 a/\nnew line after line 5\n/
Edit /ERROR/ p
\end{Verbatim}

\section{Structural Commands}
\label{sec:structural}

\index{structural commands}

The structural commands are what make sam's language truly powerful.
They apply commands to \emph{each match} of a pattern, rather than to
lines.

\begin{longtable}{@{}lp{3.5in}@{}}
\toprule
Command & Action \\
\midrule
\endhead
\texttt{x/\textit{regexp}/\textit{cmd}} & \textbf{Extract}: for each
  match of \textit{regexp} in the addressed text, set dot to the match
  and execute \textit{cmd}. \\
\texttt{y/\textit{regexp}/\textit{cmd}} & \textbf{Complement}: for
  each piece of text between matches of \textit{regexp}, set dot and
  execute \textit{cmd}. \\
\texttt{g/\textit{regexp}/\textit{cmd}} & \textbf{Guard}: execute
  \textit{cmd} only if the addressed text contains a match of
  \textit{regexp}. \\
\texttt{v/\textit{regexp}/\textit{cmd}} & \textbf{Inverse guard}:
  execute \textit{cmd} only if the addressed text does \emph{not}
  contain a match of \textit{regexp}. \\
\bottomrule
\end{longtable}

The key insight from Pike's ``Structural Regular Expressions'' paper is
that \texttt{x} replaces the traditional line-oriented approach.
Instead of ``for each line,'' you say ``for each match.''

\paragraph{Idiom: ``for each line'' is \texttt{x/.*\textbackslash n/}.}

Since \texttt{x} matches arbitrary patterns, it operates on structure
rather than lines.  You can loop over words, sentences, paragraphs,
function bodies, or any textual pattern.

\section{Grouping}
\label{sec:grouping}

\index{grouping commands}

Multiple commands can be grouped with braces:

\begin{Verbatim}[frame=single]
Edit ,x/.*\n/ {
    g/TODO/ p
}
\end{Verbatim}

\noindent
This prints every line containing ``TODO''.

\section{Multi-File Commands}
\label{sec:multifile}

\index{X command}
\index{Y command}

\begin{longtable}{@{}lp{3.5in}@{}}
\toprule
Command & Action \\
\midrule
\endhead
\texttt{X/\textit{regexp}/\textit{cmd}} & For each open file whose
  name matches \textit{regexp}, execute \textit{cmd}. \\
\texttt{Y/\textit{regexp}/\textit{cmd}} & For each open file whose
  name does \emph{not} match \textit{regexp}, execute \textit{cmd}. \\
\bottomrule
\end{longtable}

\noindent
These operate across all open windows, not just the current one.

\paragraph{Example:} Replace ``foo'' with ``bar'' in all open \file{.c}
files:

\begin{Verbatim}[frame=single]
Edit X/\.c$/ ,s/foo/bar/g
\end{Verbatim}

\section{Substitution Details}
\label{sec:subst}

\index{substitute command}

The substitute command supports backreferences:

\begin{itemize}
\item \texttt{\textbackslash 1}, \texttt{\textbackslash 2}, \ldots\
      refer to captured groups in the regexp.
\item \texttt{\&} refers to the entire match.
\end{itemize}

\paragraph{Example:} Swap the order of two-word variable names:

\begin{Verbatim}[frame=single]
Edit ,s/([a-z]+)_([a-z]+)/\2_\1/g
\end{Verbatim}

\section{Worked Examples}
\label{sec:examples}

These examples demonstrate common editing tasks using the Edit command.

\subsection{Delete All Blank Lines}

\begin{Verbatim}[frame=single]
Edit ,x/\n\n+/ c/\n/
\end{Verbatim}

\noindent
Translation: in the whole file (\texttt{,}), for each run of two or
more newlines (\texttt{x/\textbackslash n\textbackslash n+/}), replace
with a single newline (\texttt{c/\textbackslash n/}).

\subsection{Indent Selected Lines}

\begin{Verbatim}[frame=single]
Edit x/.*\n/ i/\t/
\end{Verbatim}

\noindent
For each line in the selection, insert a tab at the beginning.  (Select
the lines first, then run this Edit command.)

\subsection{Remove Trailing Whitespace}

\begin{Verbatim}[frame=single]
Edit ,s/[ \t]+$//g
\end{Verbatim}

\subsection{Rename a Variable}

\begin{Verbatim}[frame=single]
Edit ,s/\boldName\b/newName/g
\end{Verbatim}

\noindent
The \texttt{\textbackslash b} word boundaries prevent partial matches.

\subsection{Delete Lines Containing ``DEBUG''}

\begin{Verbatim}[frame=single]
Edit ,x/.*\n/ g/DEBUG/ d
\end{Verbatim}

\noindent
For each line, if it contains ``DEBUG'', delete it.

\subsection{Extract Function Signatures from a C File}

\begin{Verbatim}[frame=single]
Edit ,x/^[a-zA-Z].*\n{/ p
\end{Verbatim}

\noindent
Print every line that starts with a letter and ends with \texttt{\{},
which is a rough match for C function definitions.

\subsection{Number All Lines}

Select the text, then:

\begin{Verbatim}[frame=single]
Edit |awk '{printf "%4d  %s\n", NR, $0}'
\end{Verbatim}

\noindent
Sometimes the easiest ``Edit'' command is a pipe to a Unix tool.

\section{What Is Not Implemented}

Acme implements most of the sam command language, with these exceptions:

\begin{itemize}
\item \texttt{k} (mark) --- not implemented.
\item \texttt{n} (print file list) --- not implemented.
\item \texttt{q} (quit) --- use \cmd{Exit} instead.
\item \texttt{!} (shell command) --- use acme's native shell execution.
\end{itemize}

\section{Further Reading}

The Edit command language is fully described in:

\begin{itemize}
\item Rob Pike, ``The Text Editor sam,'' \emph{Software---Practice and
      Experience}, Vol.~17, No.~11, November 1987.
\item Rob Pike, ``Structural Regular Expressions,'' Proc.\ EUUG Spring
      1987 Conference.
\item The sam(1) manual page.
\end{itemize}

\chapter{Plumbing}
\label{ch:plumbing}

\index{plumbing}
\index{plumber}

When you right-click (\button{3}) on text in acme, the \emph{plumber}
decides what to do with it.  Is it a filename?  Open it.  A URL?
Launch a browser.  A compiler error?  Jump to the source line.  The
plumber is a pattern-matching system that routes text to the appropriate
action.

This standalone acme has an embedded plumber---no external daemon or
configuration is needed.  It comes with sensible defaults and supports
custom rules.

\section{How Plumbing Works}

When you right-click on text, acme:

\begin{enumerate}
\item Identifies the text under the click (expands to word/filename
      boundaries).
\item Constructs a \emph{plumb message} containing the text, its type,
      the source directory, and any attributes (like an address).
\item Matches the message against plumbing rules in order.
\item The first matching rule determines the action: open in the editor,
      start an external program, or search.
\end{enumerate}

\section{Default Rules}
\label{sec:defaultrules}

\index{plumbing!default rules}

The built-in rules handle common patterns:

\begin{longtable}{@{}lp{3.2in}@{}}
\toprule
Pattern & Action \\
\midrule
\endhead
URLs (\texttt{http://}, \texttt{https://}, etc.) &
  Open in web browser. \\
Image files (\texttt{.jpg}, \texttt{.png}, \texttt{.gif}, etc.) &
  Open in image viewer. \\
Document files (\texttt{.pdf}, \texttt{.ps}, \texttt{.dvi}) &
  Open in document viewer. \\
\texttt{file:line.col,line.col} &
  Open file at line and column range. \\
\texttt{file:line.col} &
  Open file at line and column. \\
\texttt{file:line} or \texttt{file:/regexp/} &
  Open file at address. \\
Plain filenames &
  Open file in acme. \\
\texttt{<file.h>} style includes &
  Search \file{/usr/include} and \file{/usr/local/include}. \\
\bottomrule
\end{longtable}

\subsection{Platform Differences}
\index{plumbing!platform}

The command used to open URLs, images, and PDFs varies by platform:

\begin{itemize}
\item \textbf{Linux:} \cmd{xdg-open}
\item \textbf{macOS:} \cmd{open}
\item \textbf{Windows (MSYS2):} \cmd{cygstart}
\end{itemize}

\section{Custom Rules}
\label{sec:customrules}

\index{plumbing!custom rules}
\index{plumbing rules file@\file{\textasciitilde/lib/plumbing}}

To override or extend the default rules, create the file
\file{\textasciitilde/lib/plumbing}.  If this file exists, it
\emph{replaces} the built-in defaults entirely.

\subsection{Rule Syntax}

A plumbing rule is a sequence of patterns and actions, separated by
blank lines.  Each rule consists of:

\begin{lstlisting}[frame=none,basicstyle=\small\ttfamily]
# Optional comment
type is text
data matches '<regexp>'
arg isfile $1
plumb to <port>
plumb start <command> $file
\end{lstlisting}

\paragraph{Pattern lines} test properties of the plumb message:

\begin{longtable}{@{}lp{3.2in}@{}}
\toprule
Pattern & Meaning \\
\midrule
\endhead
\texttt{type is text} & The message type is ``text''. \\
\texttt{data matches '\textit{re}'} & The message data matches the regex.
  Capturing groups set \texttt{\$1}, \texttt{\$2}, etc. \\
\texttt{data set \$file} & Set the data to the matched file path. \\
\texttt{arg isfile \$\textit{n}} & Check that the captured group is an
  existing file path.  Sets \texttt{\$file} if so. \\
\texttt{attr add addr=\$\textit{n}} & Add an attribute (e.g., a line
  number address). \\
\bottomrule
\end{longtable}

\paragraph{Action lines} specify what to do:

\begin{longtable}{@{}lp{3.2in}@{}}
\toprule
Action & Meaning \\
\midrule
\endhead
\texttt{plumb to edit} & Open in acme editor. \\
\texttt{plumb to web} & Open as a URL. \\
\texttt{plumb to image} & Open as an image. \\
\texttt{plumb client \$editor} & Open in the acme editor (for ``edit''
  port). \\
\texttt{plumb start \textit{cmd} \$file} & Run an external command. \\
\bottomrule
\end{longtable}

\subsection{Variables}

\begin{longtable}{@{}lp{3.5in}@{}}
\toprule
Variable & Meaning \\
\midrule
\endhead
\texttt{\$0} & The entire matched data. \\
\texttt{\$1}, \texttt{\$2}, \ldots & Captured groups from \texttt{data matches}. \\
\texttt{\$file} & The resolved file path (set by \texttt{arg isfile}). \\
\texttt{\$editor} & The editor variable (set at top of rules). \\
\bottomrule
\end{longtable}

\subsection{Example Custom Rules}

\paragraph{Open Go import paths:}
\begin{lstlisting}[frame=none,basicstyle=\small\ttfamily]
# Go import paths -> go doc
type is text
data matches '[a-zA-Z0-9./\-]+'
data matches '([a-z]+\.[a-z]+/[a-zA-Z0-9./\-]+)'
plumb to web
plumb start xdg-open https://pkg.go.dev/$1
\end{lstlisting}

\paragraph{Open Python files with line numbers from tracebacks:}
\begin{lstlisting}[frame=none,basicstyle=\small\ttfamily]
type is text
data matches 'File "([^"]+)", line ([0-9]+)'
arg isfile $1
data set $file
attr add addr=$2
plumb to edit
plumb client $editor
\end{lstlisting}

\paragraph{Open man pages:}
\begin{lstlisting}[frame=none,basicstyle=\small\ttfamily]
type is text
data matches '([a-zA-Z0-9_\-]+)\(([0-9])\)'
plumb to none
plumb start xterm -e man $2 $1
\end{lstlisting}

\section{Debugging Plumbing}

If right-clicking does not produce the expected result:

\begin{enumerate}
\item Check \file{\textasciitilde/lib/plumbing} for syntax errors.
\item Remember that custom rules \emph{replace} the defaults---include
      the standard rules in your file if you want to keep them.
\item Rules are matched in order; the first match wins.
\item File existence checks (\texttt{arg isfile}) use the source
      window's directory for relative paths.
\end{enumerate}

\chapter{Practical Workflows}
\label{ch:workflows}

\index{workflows}

This chapter demonstrates how to use acme for common programming tasks.
The goal is not to prescribe a workflow but to show how acme's simple
tools compose into effective patterns.

\section{Editing Code}
\label{sec:coding}

\subsection{Opening a Project}

Start acme with the files you want to edit:

\begin{lstlisting}[language=bash]
acme src/*.c include/*.h Makefile
\end{lstlisting}

Or start acme empty and right-click (\button{3}) on directory names to
browse, then right-click on filenames to open them.

\subsection{Navigating}

\begin{itemize}
\item Right-click on a function name to search for it in the current
      file.
\item Right-click on \file{header.h} to open it.
\item Right-click on \texttt{:42} to jump to line 42.
\item Right-click on \texttt{:/pattern/} to search by regex.
\item Use \cmd{Look} with a 2-1 chord to search for specific text.
\end{itemize}

\subsection{Multiple Columns}

Use two or three columns to keep related files visible:

\begin{itemize}
\item Left column: header files and interfaces.
\item Middle column: implementation files you are editing.
\item Right column: \file{+Errors}, \cmd{Win} shell, or test output.
\end{itemize}

Create columns with \cmd{Newcol}.  Move windows between columns by
dragging the layout button.

\section{Building and Testing}
\label{sec:building}

\subsection{Running Make}

Type \cmd{make} in any tag bar and middle-click it.  Output appears in
\file{+Errors}.

If the build fails with errors like:

\begin{Verbatim}[frame=single]
main.c:42: error: undeclared identifier 'x'
\end{Verbatim}

\noindent
right-click on that line to jump directly to \file{main.c} line~42.
Fix the error, then middle-click \cmd{make} again.

\subsection{Running Tests}

\begin{lstlisting}[language=bash]
make test
go test ./...
python -m pytest
cargo test
\end{lstlisting}

Any of these can be typed in a tag and executed.  Test output with
file:line references is clickable.

\subsection{Continuous Workflow}

A typical edit-build-fix cycle:

\begin{enumerate}
\item Edit code in a window.
\item Middle-click \cmd{Put} to save.
\item Middle-click \cmd{make} in the tag.
\item Right-click error lines to navigate to problems.
\item Fix and repeat.
\end{enumerate}

For even faster iteration, open a \cmd{Win} window and run build
commands interactively.

\section{Version Control}
\label{sec:vcs}

\index{version control}
\index{git}

\subsection{Checking Status}

Type in any tag and middle-click:

\begin{lstlisting}[language=bash]
git status
git diff
git log --oneline -20
\end{lstlisting}

\subsection{Committing}

\begin{lstlisting}[language=bash]
git add -p
git commit
\end{lstlisting}

For interactive commands like \cmd{git add -p} or \cmd{git commit}
(which opens an editor), use a \cmd{Win} window.

\subsection{Viewing Diffs}

Run \cmd{git diff} and the output appears in \file{+Errors}.
Right-click on filenames in the diff headers to open those files.

\section{Writing Prose}
\label{sec:prose}

\subsection{Proportional Fonts}

Switch to a proportional font for more comfortable reading:

\begin{enumerate}
\item Middle-click \cmd{Font} in the tag to toggle between fixed and
      proportional.
\item Or specify a font: \cmd{Font /mnt/font/DejaVuSerif/14a/font}
\end{enumerate}

\subsection{Reflowing Text}

Select a paragraph and execute:

\begin{lstlisting}[language=bash]
|fmt -w 72
\end{lstlisting}

This reflows the text to 72 columns, suitable for email or plain-text
documents.

\subsection{Spell Checking}

Select text and run:

\begin{lstlisting}[language=bash]
>aspell list
\end{lstlisting}

This sends the selection to \cmd{aspell} and shows misspelled words in
\file{+Errors}.

\section{Multi-File Search and Replace}
\label{sec:multifile}

\subsection{Using Edit X}

To replace ``oldFunc'' with ``newFunc'' in all open \file{.c} files:

\begin{Verbatim}[frame=single]
Edit X/\.c$/ ,s/oldFunc/newFunc/g
\end{Verbatim}

\subsection{Using Grep}

Run \cmd{grep -rn pattern *.c} from a tag.  The output lines are
clickable:

\begin{Verbatim}[frame=single]
src/main.c:42:    oldFunc(x, y);
src/util.c:17:    oldFunc(a, b);
\end{Verbatim}

Right-click each line to jump to the match, make the change, and move
to the next.

\section{Session Management}
\label{sec:sessions}

\index{session management}
\index{Dump}
\index{Load}

\subsection{Saving Your Workspace}

Before quitting, middle-click \cmd{Dump} in the top tag.  This saves
the complete state---all windows, their positions, fonts, and column
layout---to \file{\$HOME/acme.dump}.

\subsection{Restoring}

Next time you start acme:

\begin{lstlisting}[language=bash]
acme -l $HOME/acme.dump
\end{lstlisting}

Or start acme normally and middle-click \cmd{Load} in the top tag.

\subsection{Multiple Sessions}

Save different sessions to different files:

\begin{Verbatim}[frame=single]
Dump project-a.dump
Dump project-b.dump
\end{Verbatim}

\noindent
Then load the appropriate one:

\begin{Verbatim}[frame=single]
Load project-a.dump
\end{Verbatim}

\section{Tips}

\begin{itemize}
\item Add frequently-used commands to the top tag bar.  For example,
      type \cmd{make} or \cmd{git diff} there so it is always one
      click away.
\item Use directory windows as a file browser---right-click a directory
      path to list it.
\item Right-clicking on a word in \file{+Errors} output searches for
      it, which is useful for tracing error messages.
\item The \cmd{Zerox} command creates a second view of a file, handy for
      looking at one part while editing another.
\item \cmd{Sort} in a column tag alphabetizes windows, making it easy
      to find files.
\end{itemize}

\chapter{Customization}
\label{ch:customize}

\index{customization}

Acme deliberately has very little to configure.  There are no
configuration files, no themes, no plugin system.  This is by design:
the power comes from composing external tools, not from configuring the
editor.  Nevertheless, there are several things you can adjust.

\section{Fonts}
\label{sec:fonts}

\index{fonts}

Acme maintains two fonts: a proportional (variable-width) font and a
fixed-width font.  You can set both at startup and switch between them
per window.

\subsection{Startup Fonts}

\begin{lstlisting}[language=bash]
acme -f 'DejaVuSans/14a/font' -F 'DejaVuSansMono/14a/font'
\end{lstlisting}

The font name format is a fontconfig name followed by a size.  The
\texttt{a} suffix indicates anti-aliased rendering.

\subsection{Per-Window Fonts}

Middle-click \cmd{Font} in any window tag to toggle between the
proportional and fixed-width font.

To set a specific font for one window, type:

\begin{Verbatim}[frame=single]
Font /mnt/font/GoMono/13a/font
\end{Verbatim}

\noindent
and middle-click the whole thing.

\subsection{Available Fonts}

Acme discovers fonts through fontconfig.  Any TrueType or OpenType font
installed on your system is available.  To list installed fonts:

\begin{lstlisting}[language=bash]
fc-list : family | sort -u
\end{lstlisting}

\section{Tab Width}
\label{sec:tabwidth}

\index{tab width}

The default tab width is 4 columns.  To change it:

\begin{description}[style=nextline]
\item[Per window:]
Middle-click \cmd{Tab 8} (or any number) in the window tag.

\item[At startup:]
Set the \env{tabstop} environment variable:
\begin{lstlisting}[language=bash]
tabstop=8 acme
\end{lstlisting}
\end{description}

\section{Auto-Indentation}
\label{sec:autoindent}

\index{auto-indent}

Auto-indent is off by default.  When enabled, pressing \key{Enter}
copies the leading whitespace of the current line onto the new line.

\begin{description}[style=nextline]
\item[At startup:]
Use the \cmd{-a} flag:
\begin{lstlisting}[language=bash]
acme -a
\end{lstlisting}

\item[Per window:]
Middle-click \cmd{Indent on} or \cmd{Indent off} in the tag.

\item[Globally:]
Middle-click \cmd{Indent ON} or \cmd{Indent OFF} (uppercase).
\end{description}

\section{Plumbing Rules}
\label{sec:plumbrules}

\index{plumbing!customization}

The most significant customization you can make is writing custom
plumbing rules.  Create \file{\textasciitilde/lib/plumbing} with your
rules.  See Chapter~\ref{ch:plumbing} for the full syntax and examples.

Common customizations:

\begin{itemize}
\item Add rules for your programming language's error format.
\item Add rules for documentation URLs (e.g., Go import paths,
      Rust crate docs).
\item Add rules for your project's ticket system (e.g., clicking
      ``JIRA-1234'' opens the ticket in a browser).
\item Change the file opener command.
\end{itemize}

\section{Shell}
\label{sec:shellconfig}

\index{shell configuration}

\subsection{Command Shell}

The shell used for executing commands (middle-click on text) is
controlled by \env{acmeshell}:

\begin{lstlisting}[language=bash]
acmeshell=bash acme
\end{lstlisting}

If not set, acme defaults to \cmd{sh}.

\subsection{Win Shell}

The \cmd{Win} command uses \env{SHELL}:

\begin{lstlisting}[language=bash]
SHELL=/bin/zsh acme
\end{lstlisting}

\section{Helper Scripts}
\label{sec:helpers}

\index{helper scripts}

Since acme sets \env{winid} and \env{\%} in the environment of child
processes, you can write scripts that interact with the current window.

\subsection{Example: Format on Save}

Add to a tag:

\begin{Verbatim}[frame=single]
Fmt
\end{Verbatim}

\noindent
Where \cmd{Fmt} is a script in your PATH:

\begin{lstlisting}[language=bash]
#!/bin/sh
# Fmt: format the current file and reload
case "$%" in
    *.go) gofmt -w "$%" ;;
    *.py) black "$%" ;;
    *.c)  indent -kr "$%" ;;
    *.rs) rustfmt "$%" ;;
esac
\end{lstlisting}

After running \cmd{Fmt}, middle-click \cmd{Get} to reload the
formatted file.

\subsection{Example: Run Tests for Current File}

\begin{lstlisting}[language=bash]
#!/bin/sh
# Test: run tests related to the current file
case "$%" in
    *_test.go) go test -run . -v "$(dirname "$%")" ;;
    *.py)      python -m pytest "$%" -v ;;
    *.c)       make test ;;
esac
\end{lstlisting}

\subsection{Example: Open Documentation}

\begin{lstlisting}[language=bash]
#!/bin/sh
# Doc: open documentation for the word argument
word="$1"
case "$%" in
    *.go) xdg-open "https://pkg.go.dev/search?q=$word" ;;
    *.py) xdg-open "https://docs.python.org/3/search.html?q=$word" ;;
    *.rs) xdg-open "https://docs.rs/releases/search?query=$word" ;;
esac
\end{lstlisting}

\section{Window Size}
\label{sec:winsize}

\index{window size}

Set the initial window size with the \cmd{-W} flag:

\begin{lstlisting}[language=bash]
acme -W 1400x900
\end{lstlisting}

\section{Scrollbar Direction}
\label{sec:scrolldir}

\index{scrollbar!direction}

The \cmd{-r} flag reverses the scrollbar button mapping.  By default,
\button{1} in the scrollbar scrolls up and \button{3} scrolls down.
With \cmd{-r}, this is reversed.

\begin{lstlisting}[language=bash]
acme -r
\end{lstlisting}

\section{Startup Script}

A common pattern is to create a shell alias or script that starts acme
with your preferred settings:

\begin{lstlisting}[language=bash]
#!/bin/sh
# ~/bin/a - start acme with my preferences
export tabstop=4
export acmeshell=bash
exec acme -a \
    -f 'GoRegular/13a/font' \
    -F 'GoMono/13a/font' \
    "$@"
\end{lstlisting}

There is no \file{.acmerc} or configuration file by design.  A shell
script gives you the same effect with full transparency.


% ---- Appendices ----
\appendix
\chapter{Command Reference}
\label{app:cmdref}

\index{command reference}

All built-in commands, listed alphabetically.  Commands are executed by
middle-clicking (\button{2}) on the command name.

\begin{longtable}{@{}lp{3.6in}@{}}
\toprule
Command & Description \\
\midrule
\endhead
\cmd{Abort} &
  Debug command.  First execution prints a warning; second calls
  \texttt{abort()} to generate a core dump. \\
\cmd{Cut} &
  Delete selected text and copy to snarf buffer.  Marks window
  dirty. \\
\cmd{Del} &
  Close window.  Warns if dirty; execute again to force. \\
\cmd{Delcol} &
  Delete column and all its windows.  Warns if any are dirty. \\
\cmd{Delete} &
  Close window without dirty check. \\
\cmd{Dump} \textit{[file]} &
  Save session state.  Default: \file{\$HOME/acme.dump}. \\
\cmd{Edit} \textit{cmd} &
  Execute sam editing command on window body.  See
  Chapter~\ref{ch:edit}. \\
\cmd{Exit} &
  Quit acme.  Warns if any window is dirty; execute again to
  force. \\
\cmd{Font} \textit{[name]} &
  Toggle proportional/fixed font, or set font by name. \\
\cmd{Get} \textit{[file]} &
  Reload file from disk.  With argument, load named file. \\
\cmd{ID} &
  Print window ID number. \\
\cmd{Incl} \textit{[dirs...]} &
  List or add supplementary include directories for
  \texttt{<file.h>} resolution. \\
\cmd{Indent} \textit{[on|off|ON|OFF]} &
  Set auto-indent.  Lowercase: current window.  Uppercase: all
  windows. \\
\cmd{Kill} \textit{[names...]} &
  Send kill signal to named running commands. \\
\cmd{Load} \textit{[file]} &
  Restore session state.  Default: \file{\$HOME/acme.dump}. \\
\cmd{Look} \textit{[text]} &
  Literal text search in window body.  Uses selection or argument. \\
\cmd{New} \textit{[files...]} &
  Create new window.  With arguments, open named files. \\
\cmd{Newcol} &
  Create a new column. \\
\cmd{Paste} &
  Replace selection with snarf buffer contents. \\
\cmd{Put} \textit{[file]} &
  Write window to disk.  With argument, write to named file. \\
\cmd{Putall} &
  Write all dirty windows whose names are real files. \\
\cmd{Redo} &
  Redo last undone change. \\
\cmd{Send} &
  In Win windows: send pending input to shell.  In regular windows:
  append snarf buffer to body. \\
\cmd{Snarf} &
  Copy selected text to snarf buffer without deleting. \\
\cmd{Sort} &
  Sort windows in column alphabetically by name. \\
\cmd{Tab} \textit{[n]} &
  Print or set tab width for current window. \\
\cmd{Undo} &
  Undo last change or change group. \\
\cmd{Win} &
  Open interactive shell session in new window.  See
  Section~\ref{sec:win}. \\
\cmd{Zerox} &
  Clone window (second view of same file). \\
\bottomrule
\end{longtable}

\chapter{Edit Command Reference}
\label{app:editref}

\index{Edit command!reference}

This appendix summarizes the sam editing language available via acme's
\cmd{Edit} command.  See Chapter~\ref{ch:edit} for tutorial coverage.

\section{Addresses}

\begin{longtable}{@{}lp{3.5in}@{}}
\toprule
Address & Meaning \\
\midrule
\endhead
\texttt{.} & Current selection (dot). \\
\texttt{0} & Beginning of file. \\
\texttt{\$} & End of file. \\
\texttt{\textit{n}} & Line \textit{n} (1-based). \\
\texttt{\#\textit{n}} & Character \textit{n} (0-based). \\
\texttt{/\textit{re}/} & Next match of regex \textit{re} forward. \\
\texttt{?\textit{re}?} & Next match of regex \textit{re} backward. \\
\texttt{\textit{a},\textit{b}} & Range from \textit{a} to \textit{b}. \\
\texttt{\textit{a};\textit{b}} & Range; set dot to \textit{a} before evaluating
  \textit{b}. \\
\texttt{\textit{a}+\textit{b}} & \textit{b} relative to end of \textit{a}. \\
\texttt{\textit{a}-\textit{b}} & \textit{b} relative to start of \textit{a},
  backward. \\
\texttt{,} & Shorthand for \texttt{0,\$} (entire file). \\
\bottomrule
\end{longtable}

\section{Commands}

\begin{longtable}{@{}lp{3.5in}@{}}
\toprule
Command & Action \\
\midrule
\endhead
\texttt{a/\textit{text}/} & Append \textit{text} after addressed range. \\
\texttt{c/\textit{text}/} & Change: replace addressed range with \textit{text}. \\
\texttt{d} & Delete addressed range. \\
\texttt{i/\textit{text}/} & Insert \textit{text} before addressed range. \\
\texttt{m \textit{addr}} & Move addressed range to after \textit{addr}. \\
\texttt{p} & Print addressed range. \\
\texttt{s/\textit{re}/\textit{text}/} & Substitute first match of \textit{re}. \\
\texttt{s/\textit{re}/\textit{text}/g} & Substitute all matches. \\
\texttt{t \textit{addr}} & Copy addressed range to after \textit{addr}. \\
\texttt{=} & Print address (line number and character offset). \\
\bottomrule
\end{longtable}

\section{Structural Commands}

\begin{longtable}{@{}lp{3.5in}@{}}
\toprule
Command & Action \\
\midrule
\endhead
\texttt{x/\textit{re}/\textit{cmd}} & Extract: for each match of \textit{re}, set
  dot to match and run \textit{cmd}. \\
\texttt{y/\textit{re}/\textit{cmd}} & Complement: for each stretch between
  matches, set dot and run \textit{cmd}. \\
\texttt{g/\textit{re}/\textit{cmd}} & Guard: run \textit{cmd} if addressed text
  matches \textit{re}. \\
\texttt{v/\textit{re}/\textit{cmd}} & Inverse guard: run \textit{cmd} if addressed
  text does \emph{not} match \textit{re}. \\
\texttt{\{ \textit{cmds} \}} & Group multiple commands. \\
\bottomrule
\end{longtable}

\section{Multi-File Commands}

\begin{longtable}{@{}lp{3.5in}@{}}
\toprule
Command & Action \\
\midrule
\endhead
\texttt{X/\textit{re}/\textit{cmd}} & For each file whose name matches
  \textit{re}, run \textit{cmd}. \\
\texttt{Y/\textit{re}/\textit{cmd}} & For each file whose name does \emph{not}
  match \textit{re}, run \textit{cmd}. \\
\bottomrule
\end{longtable}

\section{Substitute Backreferences}

In the replacement text of \texttt{s} commands:

\begin{longtable}{@{}lp{3.5in}@{}}
\toprule
Token & Meaning \\
\midrule
\endhead
\texttt{\&} & Entire match. \\
\texttt{\textbackslash 1} \ldots\ \texttt{\textbackslash 9} &
  Captured subexpressions. \\
\texttt{\textbackslash n} & Literal newline. \\
\texttt{\textbackslash t} & Literal tab. \\
\texttt{\textbackslash\textbackslash} & Literal backslash. \\
\bottomrule
\end{longtable}

\chapter{Mouse Reference}
\label{app:mouseref}

\index{mouse!reference}

\section{Single-Button Actions}

\begin{longtable}{@{}llp{2.5in}@{}}
\toprule
Button & Location & Action \\
\midrule
\endhead
\button{1} & Body/Tag & Place cursor or select text. \\
\button{1} & Scrollbar & Scroll up (line at click moves to top). \\
\button{1} & Layout button & Drag to move window. \\
\button{1} double & Body/Tag & Select word. \\
\button{1} double & On bracket & Select to matching bracket. \\
\button{1} double & On quote & Select to matching quote. \\
\button{1} triple & Body/Tag & Select line. \\
\midrule
\button{2} & Body/Tag & Execute text as command. \\
\button{2} & Scrollbar & Jump scroll (proportional). \\
\midrule
\button{3} & Body/Tag & Open file, follow address, or search. \\
\button{3} & Scrollbar & Scroll down (top line moves to click). \\
\midrule
Scroll up & Body & Scroll text up. \\
Scroll down & Body & Scroll text down. \\
\bottomrule
\end{longtable}

\section{Chords}

\begin{longtable}{@{}lp{3.5in}@{}}
\toprule
Chord & Action \\
\midrule
\endhead
\chord{1}{2} & Cut: delete selection to snarf buffer. \\
\chord{1}{3} & Paste: replace selection with snarf buffer. \\
\chord{1}{2}+\button{3} & Snarf: copy selection to snarf buffer (no delete). \\
\button{2} hold, \button{1} click & 2-1 chord: pass \button{1} text as argument to
  \button{2} command. \\
\bottomrule
\end{longtable}

\section{Modifier-Click Emulation}

For mice with fewer than three buttons:

\begin{longtable}{@{}ll@{}}
\toprule
Modifier Click & Equivalent \\
\midrule
\endhead
\key{Ctrl}+\button{1} & \button{2} (execute) \\
\key{Alt}+\button{1} & \button{3} (look/open) \\
\bottomrule
\end{longtable}

\section{Button 3 Decision Table}

What happens when you right-click (\button{3}) on different kinds of
text:

\begin{longtable}{@{}lp{3.0in}@{}}
\toprule
Text & Action \\
\midrule
\endhead
\file{filename} & Open file (or jump to existing window). \\
\file{filename:N} & Open file at line N. \\
\file{filename:N.M} & Open file at line N, column M. \\
\file{filename:N.M,N.M} & Open file at line/column range. \\
\file{filename:/re/} & Open file, search for regex. \\
\texttt{:N} & Jump to line N in current file. \\
\texttt{:/re/} & Search for regex in current file. \\
\texttt{http://...} & Open URL in browser. \\
\texttt{<file.h>} & Search include directories. \\
Other text & Literal search in current body. \\
\bottomrule
\end{longtable}

\chapter{Keyboard Reference}
\label{app:kbdref}

\index{keyboard!reference}

\section{Cursor Movement}

\begin{longtable}{@{}lp{3.2in}@{}}
\toprule
Key & Action \\
\midrule
\endhead
\key{Left Arrow} & One character left. \\
\key{Right Arrow} & One character right. \\
\key{Up Arrow} & One line up (sticky column; scrolls at edge). \\
\key{Down Arrow} & One line down (sticky column; scrolls at edge). \\
\key{Ctrl+Left} & Beginning of line. \\
\key{Ctrl+Right} & End of line. \\
\key{Ctrl+A} & Beginning of line. \\
\key{Ctrl+E} & End of line. \\
\bottomrule
\end{longtable}

\section{Page and File Navigation}

\begin{longtable}{@{}lp{3.2in}@{}}
\toprule
Key & Action \\
\midrule
\endhead
\key{Page Up} & Scroll up one screen. \\
\key{Page Down} & Scroll down one screen. \\
\key{Ctrl+Up} & Scroll up one screen. \\
\key{Ctrl+Down} & Scroll down one screen. \\
\key{Home} & Cursor to first visible character. \\
\key{End} & Cursor to start of last visible line. \\
\key{Ctrl+Home} & Beginning of file. \\
\key{Ctrl+End} & End of file. \\
\key{Ctrl+Page Up} & Beginning of file. \\
\key{Ctrl+Page Down} & End of file. \\
\bottomrule
\end{longtable}

\section{Editing}

\begin{longtable}{@{}lp{3.2in}@{}}
\toprule
Key & Action \\
\midrule
\endhead
\key{Backspace} & Delete character left. \\
\key{Delete} & Delete character right. \\
\key{Ctrl+H} & Delete character left (same as Backspace). \\
\key{Ctrl+W} & Delete word left. \\
\key{Ctrl+U} & Delete to beginning of line. \\
\key{Enter} & Insert newline (auto-indent if enabled). \\
\key{Tab} & Insert literal tab. \\
\key{Escape} & Select recent typed text (cut on second press). \\
\bottomrule
\end{longtable}

\section{Composition}

\begin{longtable}{@{}lp{3.2in}@{}}
\toprule
Key & Action \\
\midrule
\endhead
\key{Alt} (press and release) & Toggle composition mode. \\
Then type sequence & Insert composed Unicode character. \\
\key{Alt} again & Cancel composition. \\
\bottomrule
\end{longtable}

See Appendix~\ref{app:unicode} for composition sequences.

\chapter{Unicode Composition Table}
\label{app:unicode}

\index{Unicode composition!table}

To compose a character, press and release \key{Alt} to enter composition
mode, then type the sequence shown.  The composed character is
inserted.  Press \key{Alt} again to cancel composition.

The sequences below are organized by category.  The full table contains
over 500 entries; the most commonly useful are listed here.

\section{Latin Accented Characters}

\begin{longtable}{@{}llll@{}}
\toprule
Sequence & Char & & Description \\
\midrule
\endhead
\texttt{'A} & \'{A} & & A acute \\
\texttt{'E} & \'{E} & & E acute \\
\texttt{'I} & \'{I} & & I acute \\
\texttt{'O} & \'{O} & & O acute \\
\texttt{'U} & \'{U} & & U acute \\
\texttt{'a} & \'{a} & & a acute \\
\texttt{'e} & \'{e} & & e acute \\
\texttt{'i} & \'{i} & & i acute \\
\texttt{'o} & \'{o} & & o acute \\
\texttt{'u} & \'{u} & & u acute \\
\midrule
\texttt{`A} & \`{A} & & A grave \\
\texttt{`E} & \`{E} & & E grave \\
\texttt{`a} & \`{a} & & a grave \\
\texttt{`e} & \`{e} & & e grave \\
\texttt{`o} & \`{o} & & o grave \\
\texttt{`u} & \`{u} & & u grave \\
\midrule
\texttt{\^{}A} & \^{A} & & A circumflex \\
\texttt{\^{}E} & \^{E} & & E circumflex \\
\texttt{\^{}a} & \^{a} & & a circumflex \\
\texttt{\^{}e} & \^{e} & & e circumflex \\
\texttt{\^{}o} & \^{o} & & o circumflex \\
\texttt{\^{}u} & \^{u} & & u circumflex \\
\midrule
\texttt{"A} & \"{A} & & A diaeresis \\
\texttt{"O} & \"{O} & & O diaeresis \\
\texttt{"U} & \"{U} & & U diaeresis \\
\texttt{"a} & \"{a} & & a diaeresis \\
\texttt{"o} & \"{o} & & o diaeresis \\
\texttt{"u} & \"{u} & & u diaeresis \\
\midrule
\texttt{\~{}N} & \~{N} & & N tilde \\
\texttt{\~{}n} & \~{n} & & n tilde \\
\texttt{\~{}A} & \~{A} & & A tilde \\
\texttt{\~{}a} & \~{a} & & a tilde \\
\texttt{\~{}O} & \~{O} & & O tilde \\
\texttt{\~{}o} & \~{o} & & o tilde \\
\midrule
\texttt{,C} & \c{C} & & C cedilla \\
\texttt{,c} & \c{c} & & c cedilla \\
\texttt{ss} & \ss & & sharp s (eszett) \\
\bottomrule
\end{longtable}

\section{Punctuation and Symbols}

\begin{longtable}{@{}llll@{}}
\toprule
Sequence & Char & & Description \\
\midrule
\endhead
\texttt{!!} & \textexclamdown & & Inverted exclamation \\
\texttt{??} & \textquestiondown & & Inverted question \\
\texttt{<<} & \guillemotleft & & Left guillemet \\
\texttt{>>} & \guillemotright & & Right guillemet \\
\texttt{oo} & \textdegree & & Degree sign \\
\texttt{+-} & $\pm$ & & Plus-minus \\
\texttt{xx} & $\times$ & & Multiplication \\
\texttt{-:} & $\div$ & & Division \\
\texttt{co} & \textcopyright & & Copyright \\
\texttt{ro} & \textregistered & & Registered \\
\bottomrule
\end{longtable}

\section{Greek Letters}

\begin{longtable}{@{}llllll@{}}
\toprule
Upper & & Lower & & & Name \\
\midrule
\endhead
\texttt{*A} $\Rightarrow$ A & & \texttt{*a} $\Rightarrow$ $\alpha$ & & & Alpha \\
\texttt{*B} $\Rightarrow$ B & & \texttt{*b} $\Rightarrow$ $\beta$ & & & Beta \\
\texttt{*G} $\Rightarrow$ $\Gamma$ & & \texttt{*g} $\Rightarrow$ $\gamma$ & & & Gamma \\
\texttt{*D} $\Rightarrow$ $\Delta$ & & \texttt{*d} $\Rightarrow$ $\delta$ & & & Delta \\
\texttt{*E} $\Rightarrow$ E & & \texttt{*e} $\Rightarrow$ $\epsilon$ & & & Epsilon \\
\texttt{*Z} $\Rightarrow$ Z & & \texttt{*z} $\Rightarrow$ $\zeta$ & & & Zeta \\
\texttt{*Y} $\Rightarrow$ H & & \texttt{*y} $\Rightarrow$ $\eta$ & & & Eta \\
\texttt{*H} $\Rightarrow$ $\Theta$ & & \texttt{*h} $\Rightarrow$ $\theta$ & & & Theta \\
\texttt{*I} $\Rightarrow$ I & & \texttt{*i} $\Rightarrow$ $\iota$ & & & Iota \\
\texttt{*K} $\Rightarrow$ K & & \texttt{*k} $\Rightarrow$ $\kappa$ & & & Kappa \\
\texttt{*L} $\Rightarrow$ $\Lambda$ & & \texttt{*l} $\Rightarrow$ $\lambda$ & & & Lambda \\
\texttt{*M} $\Rightarrow$ M & & \texttt{*m} $\Rightarrow$ $\mu$ & & & Mu \\
\texttt{*N} $\Rightarrow$ N & & \texttt{*n} $\Rightarrow$ $\nu$ & & & Nu \\
\texttt{*C} $\Rightarrow$ $\Xi$ & & \texttt{*c} $\Rightarrow$ $\xi$ & & & Xi \\
\texttt{*P} $\Rightarrow$ $\Pi$ & & \texttt{*p} $\Rightarrow$ $\pi$ & & & Pi \\
\texttt{*R} $\Rightarrow$ P & & \texttt{*r} $\Rightarrow$ $\rho$ & & & Rho \\
\texttt{*S} $\Rightarrow$ $\Sigma$ & & \texttt{*s} $\Rightarrow$ $\sigma$ & & & Sigma \\
\texttt{*T} $\Rightarrow$ T & & \texttt{*t} $\Rightarrow$ $\tau$ & & & Tau \\
\texttt{*U} $\Rightarrow$ $\Upsilon$ & & \texttt{*u} $\Rightarrow$ $\upsilon$ & & & Upsilon \\
\texttt{*F} $\Rightarrow$ $\Phi$ & & \texttt{*f} $\Rightarrow$ $\phi$ & & & Phi \\
\texttt{*X} $\Rightarrow$ X & & \texttt{*x} $\Rightarrow$ $\chi$ & & & Chi \\
\texttt{*Q} $\Rightarrow$ $\Psi$ & & \texttt{*q} $\Rightarrow$ $\psi$ & & & Psi \\
\texttt{*W} $\Rightarrow$ $\Omega$ & & \texttt{*w} $\Rightarrow$ $\omega$ & & & Omega \\
\bottomrule
\end{longtable}

\section{Mathematical Symbols}

\begin{longtable}{@{}llll@{}}
\toprule
Sequence & Char & & Description \\
\midrule
\endhead
\texttt{!=} & $\neq$ & & Not equal \\
\texttt{==} & $\equiv$ & & Identical / equivalent \\
\texttt{<=} & $\leq$ & & Less than or equal \\
\texttt{>=} & $\geq$ & & Greater than or equal \\
\texttt{<<} & $\ll$ & & Much less than \\
\texttt{>>} & $\gg$ & & Much greater than \\
\texttt{sr} & $\sqrt{}$ & & Square root \\
\texttt{if} & $\infty$ & & Infinity \\
\texttt{+-} & $\pm$ & & Plus-minus \\
\texttt{no} & $\neg$ & & Logical not \\
\texttt{an} & $\wedge$ & & Logical and \\
\texttt{or} & $\vee$ & & Logical or \\
\texttt{el} & $\in$ & & Element of \\
\texttt{mo} & $\ni$ & & Contains as member \\
\texttt{sb} & $\subset$ & & Subset \\
\texttt{sp} & $\supset$ & & Superset \\
\texttt{ca} & $\cap$ & & Intersection \\
\texttt{cu} & $\cup$ & & Union \\
\texttt{fa} & $\forall$ & & For all \\
\texttt{te} & $\exists$ & & There exists \\
\texttt{pd} & $\partial$ & & Partial derivative \\
\texttt{nb} & $\nabla$ & & Nabla / gradient \\
\texttt{es} & $\emptyset$ & & Empty set \\
\bottomrule
\end{longtable}

\section{Arrows}

\begin{longtable}{@{}llll@{}}
\toprule
Sequence & Char & & Description \\
\midrule
\endhead
\texttt{<-} & $\leftarrow$ & & Left arrow \\
\texttt{->} & $\rightarrow$ & & Right arrow \\
\texttt{ua} & $\uparrow$ & & Up arrow \\
\texttt{da} & $\downarrow$ & & Down arrow \\
\texttt{ab} & $\leftrightarrow$ & & Left-right arrow \\
\texttt{=>} & $\Rightarrow$ & & Double right arrow \\
\bottomrule
\end{longtable}

\section{Currency}

\begin{longtable}{@{}llll@{}}
\toprule
Sequence & Char & & Description \\
\midrule
\endhead
\texttt{eu} & \texteuro & & Euro sign \\
\texttt{l-} & \textsterling & & Pound sterling \\
\texttt{y=} & \textyen & & Yen \\
\texttt{c|} & \textcent & & Cent \\
\bottomrule
\end{longtable}

\bigskip
\noindent
The complete table of over 500 compositions is in the file
\file{lib/keyboard} in the acme source distribution.

\chapter{Default Plumbing Rules}
\label{app:plumbing}

\index{plumbing!default rules}

The following are the built-in plumbing rules compiled into acme.  They
are used when \file{\textasciitilde/lib/plumbing} does not exist.  If
you create a custom plumbing file, copy these rules as a starting point.

The opener command varies by platform: \cmd{xdg-open} on Linux,
\cmd{open} on macOS, \cmd{cygstart} on Windows/MSYS2.  In the rules
below, \texttt{OPENER} represents the platform-appropriate command.

\section{Variable Definitions}

\begin{lstlisting}[frame=none,basicstyle=\small\ttfamily]
addrelem = '((#?[0-9]+)|(/[A-Za-z0-9_\^]+/?)|[.$])'
addr     = :($addrelem([,;+\-]$addrelem)*)
twocolonaddr = ([0-9]+)[:.] ([0-9]+)
editor   = acme
\end{lstlisting}

The \texttt{addr} pattern matches sam-style addresses like \texttt{:42},
\texttt{:/pattern/}, \texttt{:\$}, etc.  The \texttt{twocolonaddr}
pattern matches \texttt{line:col} or \texttt{line.col} pairs.

\section{URLs}

\begin{lstlisting}[frame=none,basicstyle=\small\ttfamily]
type is text
data matches '(https?|ftp|file|gopher|mailto|news
    |nntp|telnet|wais|prospero)://...'
plumb to web
plumb start OPENER $0
\end{lstlisting}

Matches HTTP, HTTPS, FTP, and other URL schemes.  Opens in the
default web browser.

\section{Image Files}

\begin{lstlisting}[frame=none,basicstyle=\small\ttfamily]
type is text
data matches '[a-zA-Z0-9_\-./@]+'
data matches '(...)\.(jpe?g|JPE?G|gif|GIF|tiff?
    |TIFF?|ppm|bit|png|PNG)'
arg isfile  $0
plumb to image
plumb start OPENER $file
\end{lstlisting}

Opens image files with the system image viewer.

\section{PDF / PostScript / DVI Files}

\begin{lstlisting}[frame=none,basicstyle=\small\ttfamily]
type is text
data matches '[a-zA-Z0-9_\-./@]+'
data matches '(...)\.(ps|PS|eps|EPS|pdf|PDF|dvi|DVI)'
arg isfile  $0
plumb to postscript
plumb start OPENER $file
\end{lstlisting}

Opens document files with the system document viewer.

\section{File Addresses with Two-Part Line:Column}

\begin{lstlisting}[frame=none,basicstyle=\small\ttfamily]
# file:line.col,line.col
type is text
data matches '(filename_pattern)':$twocolonaddr,$twocolonaddr
arg isfile     $1
data set       $file
attr add       addr=$2-#0+#$3-#1,$4-#0+#$5-#1
plumb to edit
plumb client $editor
\end{lstlisting}

Handles patterns like \texttt{main.c:10.5,20.3} (range with
line.column on both ends).

\section{File Addresses with Line:Column}

\begin{lstlisting}[frame=none,basicstyle=\small\ttfamily]
# file:line.col
type is text
data matches '(filename_pattern)':$twocolonaddr
arg isfile     $1
data set       $file
attr add       addr=$2-#0+#$3-#1
plumb to edit
plumb client $editor
\end{lstlisting}

Handles patterns like \texttt{main.c:42.10} (single line.column).

\section{Existing Files with Optional Address}

\begin{lstlisting}[frame=none,basicstyle=\small\ttfamily]
type is text
data matches '(filename_pattern)'($addr)?
arg isfile  $1
data set    $file
attr add    addr=$3
plumb to edit
plumb client $editor
\end{lstlisting}

The catch-all file rule.  Matches plain filenames and filenames with
sam-style address suffixes (\texttt{main.c:42}, \texttt{main.c:/re/}).

\section{Include Files}

\begin{lstlisting}[frame=none,basicstyle=\small\ttfamily]
# .h files in /usr/include
type is text
data matches '([a-zA-Z0-9/_\-]+\.h)'($addr)?
arg isfile  /usr/include/$1
data set    $file
attr add    addr=$3
plumb to edit
plumb client $editor

# .h files in /usr/local/include
type is text
data matches '([a-zA-Z0-9/_\-]+\.h)'($addr)?
arg isfile  /usr/local/include/$1
data set    $file
attr add    addr=$3
plumb to edit
plumb client $editor
\end{lstlisting}

Two rules search standard system include directories.  Additional
directories can be added with the \cmd{Incl} command or in custom
plumbing rules.


% ---- Back Matter ----
\backmatter
\chapter*{Bibliography}
\addcontentsline{toc}{chapter}{Bibliography}
\label{ch:biblio}

\begin{description}[style=nextline,leftmargin=0pt,labelindent=0pt]

\item[Pike, Rob. ``Acme: A User Interface for Programmers.'']
Proceedings of the Winter 1994 USENIX Conference, San Francisco, 1994.\\
\url{https://research.swtch.com/acme.pdf}\\
The foundational paper describing acme's design, philosophy, and
interface.  Essential reading for understanding why acme works the way
it does.

\item[Pike, Rob. ``The Text Editor sam.'']
\emph{Software---Practice and Experience}, Vol.~17, No.~11,
pp.~813--845, November 1987.\\
\url{https://doc.cat-v.org/plan_9/4th_edition/papers/sam/}\\
The full paper on sam's design and command language, which acme's
\cmd{Edit} command implements.

\item[Pike, Rob. ``Structural Regular Expressions.'']
Proceedings of the EUUG Spring 1987 Conference, Helsinki, 1987.\\
\url{https://doc.cat-v.org/bell_labs/structural_regexps/}\\
Describes the theoretical foundation for sam's \texttt{x} and \texttt{y}
commands---applying regular expressions to structure rather than lines.

\item[Cox, Russ. ``A Tour of Acme'' (screencast).]
\url{https://research.swtch.com/acme}\\
\url{https://www.youtube.com/watch?v=dP1xVpMPn8M}\\
A video demonstration of day-to-day acme use by one of its primary
maintainers.  The best introduction for new users.

\item[acme(1) manual page.]
\url{https://9fans.github.io/plan9port/man/man1/acme.html}\\
The official manual page from plan9port.

\item[sam(1) manual page.]
\url{https://9fans.github.io/plan9port/man/man1/sam.html}\\
Reference for the sam command language used by acme's \cmd{Edit}
command.

\item[Pike, Rob. ``Plumbing and Other Utilities.'']
\url{https://doc.cat-v.org/plan_9/4th_edition/papers/plumb/}\\
Design and implementation of the plumber, Plan~9's inter-application
message routing system.

\item[Pike, Rob. ``A Minimalist Global User Interface.'']
\url{https://doc.cat-v.org/plan_9/4th_edition/papers/rio/}\\
Describes rio, Plan~9's window system, which shares acme's philosophy of
text and files as the universal interface.

\item[Plan 9 from User Space (plan9port).]
\url{https://github.com/9fans/plan9port}\\
The port of Plan~9 user-space tools to Unix, from which this standalone
acme is derived.

\item[Plan 9 Fourth Edition Papers.]
\url{https://doc.cat-v.org/plan_9/4th_edition/papers/}\\
The complete collection of Plan~9 documentation.

\item[Acme community resources.]
\url{https://acme.cat-v.org/}\\
Community-maintained collection of acme tips, scripts, and links.

\end{description}


\printindex

\end{document}
